% PDF 342--400；原书第 315--373 页。
\chapter{生产环境中的大语言模型}
\label{ch:production}
\chaptercredit{8}
\chapterabstract{大语言模型的潜力主要由研究工作推动。在研究中，任务表现常常比实际使用的其他方面受到更多关注。将这种潜力转化为能在生产环境中使用的应用，正在成为新的研究前沿。推动这一变化的有希望在市场中形成差异、优化运营或通过应用大语言模型创造独特价值的企业。本章介绍实际应用大语言模型的若干重要方面，帮助读者把模型能力转化为可用的应用。从选择大语言模型到优化模型延迟，本章重点介绍有关工具和技术，为读者开发大语言模型应用提供指导。}

% PDF 342--343；原书第 315--316 页。
\section{引言}
本章汇总开发人员构建生产环境中的大语言模型应用时应考虑的不同因素，希望用当时的实践经验和知识，帮助读者开发可靠、成本合理且安全的应用。如前所述，大语言模型既有很大的潜力，也带来风险。开发人员应当能够处理开发周期的各个步骤，尽可能发挥模型的作用，同时尽量减少风险。

第~\ref{sec:ch08-applications}~节从大语言模型的常见应用讲起，说明后续各节所涉及的使用场景。还将考察大语言模型的几种主要类别，为读者评估模型能否适应不同场景提供另一个维度。模型上下文长度、参数数量和架构等细节已在其他地方详细讨论，例如第~\ref{ch:pretraining}~章，本章不再展开。

第~8.3 和 8.4~节介绍评估大语言模型应用的常用指标，并列出多种常用于评估不同应用的代表性数据集。

第~8.5~节从开源与闭源的角度讨论大语言模型选型，考察任务表现、成本、数据安全及许可等因素，使读者了解设计应用时可能需要作哪些权衡。也会讨论推理延迟和大语言模型定制，说明选用开源或闭源模型可能给项目带来哪些限制。

第~8.6~节介绍当时迅速演变的大语言模型应用开发环境中的工具、框架和常用模式，包括可用的开发框架、提示工程工具、向量存储和大语言模型定制。

随后，第~8.7~节深入讨论推理，涉及模型托管方案、改善性能的方法，以及尤为重要的成本优化。如第~\ref{ch:adaptation}~章第~\ref{sec:ch04-compute-efficient}~节所述，大语言模型的推理开销仍是研究的重点，因此有必要了解当时的优化方法。

本章最后从大语言模型运维（LLMOps）的角度概述应用开发。大语言模型较复杂，在应用中的使用还处于初期，项目需要有明确的开发和运维流程。模型本身、交互方式和定制方法都有较大的变化空间，适当的流程使开发团队能够持续集成新方法并开展试验、评估和部署，尽量减少对用户的干扰。在大语言模型应用的早期，维持用户信任和应用的可信度很重要，LLMOps 有助于这一过程。

% PDF 343--349；原书第 316--322 页。
\section{大语言模型的应用}
\label{sec:ch08-applications}
讨论开发可用于生产环境的大语言模型应用的技术细节之前，先了解这类模型已用于哪些问题和场景。这里简要介绍大语言模型能够改善结果的几种通用应用，例如对话式聊天机器人，再概述可用于这些场景的不同类型的模型。

这些应用可以帮助读者把本章较技术性的内容放回开发周期中理解，看到它们在实践中的用途。

\subsection{对话式人工智能、聊天机器人和人工智能助手}
这是大语言模型带来显著改进的最常见场景。在聊天机器人和对话式人工智能中，与传统语言模型相比，大语言模型有更强的语言理解能力，带来了几项重要的新能力，其中尤为重要的是自然语言理解（natural language understanding，NLU）（\refcite{ch08-ref-098}{Wei et al., 2022}）。在这类应用中，大语言模型能否理解用户查询背后的意图，并将输入与已有参数知识综合起来，生成连贯的回答，对应用能否发挥作用影响很大。

用户通常进行开放式对话，可能涉及多个知识领域或主题，因此大语言模型能否跟踪上下文，也关系到整个对话会话中的回答是否连贯。在用户与应用来回交流的多轮对话中，大语言模型需要有选择地纳入前面的查询或回答，从而始终给出有用、连贯的回答。输入上下文长度的提升进一步增强了这一能力（\refcite{ch08-ref-069}{Pawar et al., 2024}），因为模型可以在更长的输入范围上进行推理。

\begin{keypoint}
这类大语言模型应用中的许多场景，对回答的事实性有较高要求。开发时需要尽量减少模型容易产生的幻觉。常见办法是使用外部知识库，从中提取相关上下文，以经过核查的知识作为模型回答的条件。怎样把知识库纳入应用架构，并把其中的知识放入模型输入以得到合适的回答，也给应用开发带来一组问题。
\end{keypoint}

\subsection{内容创作}
在内容营销和数字媒体领域，有一句常说的话叫“内容为王”。以往，故事、博客文章、通讯、社交媒体内容等，主要由有经验的人创作。他们知道哪些内容能引起受众共鸣，制作这些内容，并通过合适的渠道传播。当时已有不少内容制作环节可由大语言模型完成。有些应用允许营销人员整理有关受众人口特征的背景资料，提供相关内容和详细指导，由含有大语言模型的系统生成有吸引力的内容，并按预定日程在各渠道发布。

同样，可以用提示让大语言模型生成整篇文章，或创作关于真实或虚构主题与事件的故事。读者常常难以将这些内容与人创作的内容区分开来，内容创作者可以借此扩大创作规模，增加内容的多样性。这些发展也带来负面后果，在教育中尤其明显。学生很快开始用 OpenAI 的 ChatGPT 等大语言模型完成作业，可能未能充分掌握知识（\refcite{ch08-ref-053}{Lo, 2023}）。另一方面，大语言模型大幅提高了教育工作者设计、规划和制作课程内容的效率。使用得当时，它也是方便的学习辅助工具。

\subsection{搜索、信息检索和推荐系统}
这类应用主要是尽可能高效地为用户提供相关信息。检索常常像大海捞针。如果没有识别并呈现相关内容的办法，用户就会不断收到不切题或无关的信息。以一家公司为例，它有 10 年的产品文档，涉及多个产品版本及其功能，在这样的知识库中查找信息会很复杂。借助大语言模型，用户可以与人工智能助手对话。助手使用整合用户查询、知识库及其元数据的技术，使用户更高效地了解产品功能。

大语言模型的自然语言理解、实体抽取、摘要和语义嵌入能力有所改善，因此能以更好的精确率和召回率找出相关知识，进行推理，并将结果纳入个性化回答。这些变化有赖于大语言模型在方法、架构和工具上的许多改进，例如第~\ref{ch:rag}~章介绍的检索增强生成技术。

\urldef\githuturl\url{https://madnight.github.io/githut/#/}
\subsection{编程}
大语言模型不仅善于生成自然语言，也能生成计算机编程语言。原书以当时常用的 GitHub Copilot 为例。它为人类程序员开发软件提供代码辅助，下面用其主要功能说明此类大语言模型应用可以提供的帮助。

\begin{itemize}
\item \textbf{代码自动补全。}可以简单到像传统 Tab 补全那样补全函数、方法和变量名，也可以根据对当前脚本已有代码的实时分析，建议完整的代码块。
\item \textbf{支持多种编程语言。}使开发人员能更高效地在多种语言间协同工作。这在全栈或专业领域应用开发中尤其有用，因为方案的不同组件可能用不同语言编写。例如，一名全栈开发人员用 JavaScript 为用户界面编写数据处理程序时，Copilot 可以用 Python 建议后端 API 的代码块，为前端提供数据。据原书当时的描述，GitHub Copilot 支持公开 GitHub 仓库中可用的编程语言，其在不同语言上的能力取决于这些语言在公开仓库代码中的比重。GitHut\footnote{\githuturl} 展示了不同编程语言在 GitHub 上的相对比重。
\item \textbf{自然语言理解。}用户可以用自然语言说明希望代码实现哪些功能，往往只需编写描述下一段代码用途的注释。如果当前脚本启用了 GitHub Copilot，它会推荐实现注释所述功能的代码。对许多问题，这类建议可以给出简洁的解法，并能利用正在开发的脚本的其他上下文，尤其在注释和文档写得较好的情况下。这个功能能提高效率。与其他大语言模型生成内容一样，用户仍需核查并测试推荐的代码，防止因模型生成错误而出问题。
\item \textbf{代码重构与调试。}GitHub Copilot 的另一个用途是提高重构和调试的效率。训练使用了大量代码，模型学到了同一问题或相似问题的多种解法，因此可以提供其他模式供用户考虑。对于定义 Python 类等重复性工作，Copilot 可以生成样板代码，使开发者把时间用于具体功能。在调试时，GitHub Copilot 可以结合代码和开发人员对预期功能的描述，解读执行错误或意外输出，帮助定位可能出问题的代码。它也可以用自然语言说明执行流程和代码功能，协助开发人员查找执行问题的成因。
\end{itemize}

这些能力有助于提高软件开发效率。以往，要写出可以完成预期功能的代码，可能需要查阅参考书，多次访问 Stack Overflow 或 GitHub Gists 等网站，并请同事审查代码。现在，编程助手让开发人员可以通过同一个界面获得类似的学习和反馈。编程助手也已经集成进 Visual Studio Code、Vim 和 JetBrains 等常用集成开发环境。%
\footnote{译注：原书将 Visual Studio Code 写成“Visual Basic Code”。}

\subsection{大语言模型的类别}
\label{sec:ch08-categories}
大语言模型的数量迅速增加，开源领域尤其如此（\refcite{ch08-ref-027}{Gao and Gao, 2023}）。训练数据、微调方法、提示、算力规模、学习目标和其他影响开发的因素都可以调整，由此产生许多改进，推动模型多种能力的提升。单靠一章很难综述这些变化。不过，有许多在线资源可以帮助读者辨别具体模型在任务与基准上的表现。例如，Hugging Face 的 \texttt{open\_llm\_leaderboard} 随新模型开发与发布汇总大语言模型的任务表现。%
\footnote{\url{https://huggingface.co/spaces/HuggingFaceH4/open_llm_leaderboard}}
这类工具常允许用户按关心的因素筛选，例如只比较采用某种架构或某种权重精度的大语言模型。

比较各个模型是为具体场景选型的重要步骤，了解大语言模型的主要类别也很重要。不同类别擅长的应用不同，例如对话、编程与检索。下面介绍几种主要类别，帮助读者作出选择。

\subsubsection{通用大语言模型}
通用大语言模型通常在来自互联网的大型语料库上训练。一些语料经过筛选，但整体上，这类训练集广泛涵盖互联网上的不同主题和领域。因此，模型擅长通用语言任务，例如自然语言理解、实体抽取、思维链与问答，但未必擅长特定领域的任务，例如金融数值推理。这类大语言模型用途最广，可以作为许多应用的起点。若模型在某些方面能力不足，不少模型还可以针对任务做专门化处理，不过开源模型通常比闭源模型更便于这样做。概括地说，可以从文本生成、自然语言理解、问答、翻译、分类、摘要和对话等方面考察这类模型。对多数用例，它们直接使用时的表现可能已足够。

\subsubsection{多模态大语言模型}
多模态大语言模型在不止一种“模态”（modality）的数据上训练。常见模态包括文本、音频、视频和图像（\refcite{ch08-ref-105}{Yin et al., 2024}）。在这些模态上训练模型，带来一组新的跨模态应用，正在成为生成式人工智能的新前沿（详见第~9~章）。第~\ref{sec:ch08-applications}~节讨论的大语言模型应用，还可以借助多模态大语言模型扩展，例如依据自然语言描述检索图像，或根据自然语言指令生成音频。将文本与图像生成结合，用来讲故事或编写产品说明等，可以提供更丰富、更具表现力的用户体验，正在受到模型开发者关注。总体而言，与传统大语言模型相比，这些模型规模很大，给采用和集成带来新的困难。研究者正在继续改善这类模型的任务表现和计算效率。

\subsubsection{多语言大语言模型}
多语言大语言模型在不止一种自然语言的文本数据上训练。它们已受到广泛研究，可用于翻译、多语言推理、多语言内容生成等任务。有些模型支持的语言种类很多。例如，BigScience 开发的开源模型 BLOOM（\refcite{ch08-ref-103}{Workshop et al., 2023}），按原书此处的描述，具有 1960 亿参数，在 46 种自然语言和 13 种编程语言的文本上训练。%
\footnote{译注：本书引用的 BigScience 论文题名与第~4.4.1.1~节所用的 BLOOM-176B 均指 1760 亿参数模型。}
这类模型适合跨语言任务，例如对多种语言的文本作推理（如 \refcite{ch08-ref-076}{Ranaldi et al., 2024}）。在应用中，多语言大语言模型已用于客服等沟通场景。这些场景常出现语码转换（code-switching），即在交际中交替使用不同自然语言（\refcite{ch08-ref-106}{Yong et al., 2023}）。

\subsubsection{领域专用大语言模型}
领域专用大语言模型在来自特定行业、领域或专业方向的数据上训练，数据经过有针对性的选择。训练的目的通常是让模型的能力适应领域的特点。当领域有许多难以按通用语境理解的行话，或者其内容分布与通用语境差异较大时，这一点尤其重要。例如，在生物医学领域，研究者开发了领域专用大语言模型 BioMistral（\refcite{ch08-ref-040}{Labrak et al., 2024}）。他们在网上大型生物医学研究文献库之一 PubMed Central 上，对 Mistral 模型做适应性预训练。由 70 亿参数的通用 Mistral 适配得到的 BioMistral 模型，在 10 项生物医学任务中的 9 项优于通用模型。教育、法律、经济、政治、科学和金融等领域也有专用大语言模型，可以作为开发许多领域应用的起点。

% PDF 349--356；原书第 322--329 页。
\section{大语言模型评估指标}
\label{sec:ch08-metrics}
评估大语言模型，是系统衡量它们能否有效完成特定任务的过程。大语言模型的评估指标可按应用与方法，从不同维度分类。下面介绍其中几种。

\begin{enumerate}
\item \textbf{有参考与无参考}
  \begin{itemize}
  \item \textbf{有参考。}将模型输出与预定义的正确答案作比较。常用于翻译和摘要等任务。
  \item \textbf{无参考。}根据模型输出的内部一致性与语言性质评估质量。
  \end{itemize}
\item \textbf{基于字符、基于词语与基于嵌入}
  \begin{itemize}
  \item \textbf{基于字符。}侧重字符层面的准确性，适用于某些文本生成任务。
  \item \textbf{基于词语。}考察词语是否出现、出现频率及顺序。
  \item \textbf{基于嵌入。}借助向量表示评估语义相似度，不限于词语的精确匹配。
  \end{itemize}
\item \textbf{通用指标与任务专用指标}
  \begin{itemize}
  \item \textbf{通用指标。}无需修改即可用于不同类型的任务。
  \item \textbf{任务专用指标。}为特定应用设计，直接关系到任务输出的质量。
  \end{itemize}
\item \textbf{人工评估与大语言模型评估}
  \begin{itemize}
  \item \textbf{人工评估。}由人判断模型输出的质量或相关性。人工评估有主观性，但对判断自然语言是否流畅、前后连贯，以及在具体使用情境中是否切题十分重要。人工评估能看出自动指标可能忽略的方面，尤其是文本是否符合上下文，以及对人类语言细微含义的理解。
  \item \textbf{大语言模型评估。}用第二个大语言模型评估模型输出，通常采用自动指标或让模型作判断。
  \end{itemize}
\item \textbf{传统指标与非传统指标}
  \begin{itemize}
  \item \textbf{传统指标。}关注模型输出在词汇与句法上的准确程度。常见指标包括字符串精确匹配、字符串编辑距离、BLEU 和 ROUGE，侧重词语与短语的顺序及准确性。
  \item \textbf{非传统指标。}利用语言模型的能力，从更多方面评估生成文本的质量。例如，BERTScore 用嵌入比较语义相似度，G-Eval 等方法用另一个能力较强的大语言模型评估生成文本的质量。
  \end{itemize}
\end{enumerate}

下面介绍一些常用于大语言模型评估的指标，说明它们的计算方法和具体应用。

\subsection{困惑度}
困惑度（perplexity）衡量模型在多大程度上能预测语料中的词元序列。困惑度越低，表明模型越能准确预测序列，实际数据越符合它的预期。困惑度较高，则表明根据模型生成的下一词元概率，这个序列较为出乎模型预期。

给定分词后的序列 $X=(x_0,x_1,\ldots,x_N)$，原书以 $N$ 表示词元数。$X$ 的困惑度按下式计算。%
\footnote{译注：原书同时使用 $N$ 作词元数，并将求和指标写为 $0$ 至 $N$，共 $N+1$ 项。计算时分母应与实际计入的预测项数一致。}
\begin{equation}
\operatorname{PPL}(X)=\exp\!\left\{-\frac{1}{N}
 \sum_{i=0}^{N}\log p_\theta(x_i\mid x_{<i})\right\}.
\label{eq:ch08-perplexity}
\end{equation}
其中，$\log p_\theta(x_i\mid x_{<i})$ 表示模型以所有前置词元 $x_{<i}$ 为条件，给第 $i$ 个词元分配的对数似然。它反映模型对序列中各个词元的预测准确程度。

\subsection{BLEU}
这一类指标中常用的是 \refcite{ch08-ref-067}{Papineni 等人（2002）}提出的双语评估替补分数（Bilingual Evaluation Understudy，BLEU），主要用于评估从一种自然语言翻译成另一种自然语言的文本质量。BLEU 考察机器生成文本与一份或多份参考译文中连续词语序列的出现情况及频率。这种连续序列称为 $n$ 元组。由此衡量两者的接近程度。

BLEU 的计算包括以下几步。

\begin{itemize}
\item \textbf{计算精确率。}计算不同长度的 $n$ 元组的精确率。对于长度为 $n$ 的 $n$ 元组，精确率 $p_n$ 等于生成文本中与参考文本匹配的 $n$ 元组数，除以生成文本中的 $n$ 元组总数。每种 $n$ 元组的匹配计数，以它在任一参考文本中的最大出现次数为上限，避免过度计数。
\begin{equation}
p_n=\frac{\text{截断后的匹配 }n\text{ 元组数}}
          {\text{生成文本中的 }n\text{ 元组总数}}.
\label{eq:ch08-bleu-precision}
\end{equation}
\item \textbf{精确率的几何平均。}算出不同 $n$ 元组长度的精确率后，取所考虑长度上各精确率的几何平均，得到记作 BLEU-N 的分数。
\begin{equation}
\mathrm{BLEU\text{-}N}
=\left(\prod_{n=1}^{N}p_n\right)^{1/N}.
\label{eq:ch08-bleu-geomean}
\end{equation}
\item \textbf{长度惩罚。}较短的生成文本出错的机会较少，可能更容易得到较高的精确率。为减轻对短文本的偏好，BLEU 加入长度惩罚（brevity penalty）。生成文本长度 $c$ 短于参考长度 $r$ 时受到惩罚，以避免译文过度简短。
\begin{equation}
\mathrm{BP}=
\begin{cases}
1,& c>r,\\
e^{(1-r/c)},& c\leq r.
\end{cases}
\label{eq:ch08-bleu-penalty}
\end{equation}
\item \textbf{最终 BLEU 分数。}把各精确率的几何平均值与长度惩罚结合，得到 0 至 1 之间的分数。1 表示与参考文本完全匹配。
\begin{equation}
\mathrm{BLEU}
=\mathrm{BP}\cdot\exp\!\left(\sum_{n=1}^{N}\frac{1}{N}\log p_n\right).
\label{eq:ch08-bleu}
\end{equation}
\end{itemize}

BLEU 虽然应用广泛，但不考虑生成文本在语义上是否准确、语法上是否正确。因此，它可用于初步评估翻译质量。为更全面地反映语言生成的细微差别，还应辅以其他指标或人工评估。

\subsection{ROUGE}
面向召回率的摘要评估替补指标（Recall-Oriented Understudy for Gisting Evaluation，ROUGE）侧重召回率。它主要考察参考文本中有多少词也出现在机器生成的输出中，因此尤其适合评估自动摘要任务。

ROUGE 有多种变体，各有侧重。

\begin{itemize}
\item \textbf{ROUGE-N。}ROUGE 包括一组指标，用于比较生成文本与一组人工参考摘要，以评估摘要和翻译效果。由于侧重召回率，ROUGE 主要衡量生成文本再现参考摘要中词语和短语的程度。在需要尽量涵盖参考内容时，这一特点尤其有用。
\begin{equation}
\mathrm{ROUGE\text{-}N}=
\frac{\sum_{S\in\{\text{参考摘要}\}}\sum_{\mathrm{gram}_n\in S}
       \operatorname{Count}_{\mathrm{match}}(\mathrm{gram}_n)}
     {\sum_{S\in\{\text{参考摘要}\}}\sum_{\mathrm{gram}_n\in S}
       \operatorname{Count}(\mathrm{gram}_n)}.
\label{eq:ch08-rouge}
\end{equation}
式中，$\mathrm{gram}_n$ 表示长度为 $n$ 的 $n$ 元组，$\operatorname{Count}_{\mathrm{match}}(\mathrm{gram}_n)$ 是某种 $n$ 元组在候选摘要与参考摘要集合中共同出现的最大次数。例如，ROUGE-1 使用一元组，ROUGE-2 使用二元组。
\item \textbf{ROUGE-L。}关注生成文本与参考文本的最长公共子序列（longest common subsequence，LCS）。与 $n$ 元组重叠不同，LCS 不要求序列连续，因而能捕捉更灵活的匹配。
\item \textbf{ROUGE-W。}ROUGE-L 的扩展，在评估中考虑文本长度，以抵消长度偏差。
\item \textbf{ROUGE-S。}度量跳跃二元组的共现，即按句子顺序取任意两个词，无论中间间隔多少词。该指标侧重内容出现的顺序，允许中间插入其他内容。
\item \textbf{ROUGE-SU。}在 ROUGE-S 的基础上，同时使用跳跃二元组和一元组。
\end{itemize}

\subsection{BERTScore}
\refcite{ch08-ref-111}{Zhang 等人（2020）}提出的 BERTScore 利用 BERT 模型的深层上下文嵌入来评估文本质量。BLEU 和 ROUGE 等传统指标考察词元或 $n$ 元组的重叠，BERTScore 则用这些上下文嵌入，计算候选文本各词元与参考文本各词元的相似度。

BERTScore 用贪心匹配，使候选文本的每个词元与参考文本中最相似的词元配对，以得到较高的整体相似度。评估包括三项指标。

\begin{itemize}
\item \textbf{召回率（$R_{\mathrm{BERT}}$）。}对参考文本的每个词元取最大的相似度分数，再求和并除以参考文本中的词元数。它反映候选文本涵盖了多少参考文本内容。
\begin{equation}
R_{\mathrm{BERT}}=\frac{1}{|x|}\sum_{x_i\in x}
 \max_{\widehat{x}_j\in\widehat{x}}\langle x_i,\widehat{x}_j\rangle.
\label{eq:ch08-bert-recall}
\end{equation}
\item \textbf{精确率（$P_{\mathrm{BERT}}$）。}类似地，对候选文本中的每个词元取最大的相似度分数，求和后除以候选文本中的词元数。它衡量候选文本中的词元在参考文本中得到体现的程度。
\begin{equation}
P_{\mathrm{BERT}}=\frac{1}{|\widehat{x}|}\sum_{\widehat{x}_j\in\widehat{x}}
 \max_{x_i\in x}\langle\widehat{x}_j,x_i\rangle.
\label{eq:ch08-bert-precision}
\end{equation}
\item \textbf{$F_1$ 分数（$F_{\mathrm{BERT}}$）。}精确率与召回率的调和平均，兼顾完整性和精确程度。
\begin{equation}
F_{\mathrm{BERT}}=2\frac{P_{\mathrm{BERT}}\cdot R_{\mathrm{BERT}}}
 {P_{\mathrm{BERT}}+R_{\mathrm{BERT}}}.
\label{eq:ch08-bert-f1}
\end{equation}
\end{itemize}

BERTScore 考虑语义，对改写有一定鲁棒性，适合评估翻译或摘要。不过，它需要较多算力，结果也未必总与人工判断一致，尤其在评估文本结构和连贯性时。

\subsection{MoverScore}
MoverScore 利用词移动距离（Word Mover's Distance，WMD）的概念，评估系统预测文本与参考文本的语义相似度（\refcite{ch08-ref-038}{Kusner et al., 2015}）。该指标有助于衡量词语和短语之间的语义距离，适用于文本评估。BERTScore 为词元选择单个匹配，即“硬对齐”。MoverScore 则使用多对一匹配，即“软对齐”，使词元间的匹配更灵活。

MoverScore 的主要组成部分如下。

\begin{itemize}
\item \textbf{运输流量矩阵（$F$）。}元素 $F_{ij}$ 表示从预测序列 $x^n$ 的第 $i$ 个 $n$ 元组 $x_i^n$ 流向参考序列 $y^n$ 的第 $j$ 个 $n$ 元组 $y_j^n$ 的流量。
\item \textbf{成本矩阵（$C$）。}运输成本矩阵的每个元素 $C_{ij}$ 表示预测文本第 $i$ 个 $n$ 元组与参考文本第 $j$ 个 $n$ 元组之间的距离 $d(x_i^n,y_j^n)$。
\item \textbf{矩阵逐元素运算。}总运输成本为矩阵 $C\odot F$ 的各元素之和，其中 $\odot$ 表示逐元素相乘。
\item \textbf{$n$ 元组序列。}把系统预测文本 $x$ 和参考文本 $y$ 看成词序列，计算时使用各自的 $n$ 元组，如一元组和二元组。
\item \textbf{权重向量（$f_{x^n}$ 和 $f_{y^n}$）。}它们分别对应 $x^n$ 与 $y^n$ 中的 $n$ 元组，在各个 $n$ 元组上形成分布，通常归一化，使权重之和为 1。
\end{itemize}

WMD 是满足下列条件的运输成本最小值。
\begin{equation}
\begin{aligned}
\operatorname{WMD}(\symbf{x}^n,\symbf{y}^n)
&:=\min_{F\in\mathbb{R}^{|x^n|\times|y^n|}}\langle C,F\rangle,\\
&\phantom{:}\text{约束为 }F\symbf{1}=f_{x^n},\quad F^{\mathsf T}\symbf{1}=f_{y^n}.
\end{aligned}
\label{eq:ch08-wmd}
\end{equation}

实际计算中，MoverScore 用 $n$ 元组嵌入之间的欧氏距离衡量语义距离。原书给出下式。
\begin{equation}
d(x_i^n,y_j^n)=\lVert E(x_i^n)-E(y_j^n)\rVert_2^2.
\label{eq:ch08-mover-distance}
\end{equation}
其中，$E$ 是将 $n$ 元组映射为向量的嵌入函数。可以使用 word2vec 等传统方法，但 ELMo 和 BERT 等上下文嵌入能纳入句子层面的上下文，因此更常用。

原书所列的 $n$ 元组嵌入计算式为
\begin{equation}
E(x_i^n)=\sum_{k=i}^{i+n-1}\operatorname{idf}(x_k),
\label{eq:ch08-mover-embedding}
\end{equation}
其中，$\operatorname{idf}(x_k)$ 是 $x_k$ 的逆文档频率。%
\footnote{译注：式~\eqref{eq:ch08-mover-embedding} 右侧只有标量权重，缺少词向量。MoverScore 论文式（1）的每项为 $\operatorname{idf}(x_k)E(x_k)$。论文中的距离采用 $\ell_2$ 范数，未取平方。见 \url{https://aclanthology.org/D19-1053/}。}
每个 $n$ 元组的权重 $f_{x_i^n}$ 由下式确定。
\begin{equation}
f_{x_i^n}=\frac{1}{Z}\sum_{k=i}^{i+n-1}\operatorname{idf}(x_k),
\label{eq:ch08-mover-weight}
\end{equation}
其中，$Z$ 是归一化常数，使 $\sum f_{x^n}=1$。当 $n$ 大于句子的长度、只剩一个 $n$ 元组时，MoverScore 简化为句子移动距离（Sentence Mover's Distance，SMD）。
\begin{equation}
\operatorname{SMD}(\symbf{x}^n,\symbf{y}^n)
:=\lVert E(x_1^{l_x})-E(y^{l_y})_1\rVert,
\label{eq:ch08-smd}
\end{equation}
其中，$l_x$ 和 $l_y$ 分别为两个句子的长度。%
\footnote{译注：原书式~\eqref{eq:ch08-smd} 的后一项下标 1 印在括号外。MoverScore 论文将后一项写为 $E(y_1^{l_y})$。}

\subsection{G-Eval}
G-EVAL 用结构化、可随任务调整的方法评估生成文本，希望比传统方法更细致地衡量文本质量。它改进分数的计算与呈现方式，以应对评分差异及评估结果与人工判断的一致程度等问题。

\sourcefigure{fig08_01.pdf}{G-EVAL 框架。先将任务说明与评估准则交给大语言模型，由模型生成列出详细评估步骤的思维链。再结合思维链和初始提示，用填写评分表的方式评估自然语言生成输出。最后，按各分值的概率计算加权和。图示的连贯性加权分为 2.59。}{fig:ch08-geval}

如图~\ref{fig:ch08-geval} 所示，G-EVAL 的方法由三个主要部分组成。

\begin{enumerate}
\item \textbf{提示。}提示给出任务定义与具体评估准则。例如，评估文本摘要时，提示定义任务，指定连贯性等准则，请评分者从这些方面评价摘要。
\item \textbf{思维链。}由大语言模型生成的一系列中间指令，给出执行评估的详细步骤。思维链指导评估过程，提高评估的一致性和细致程度。例如，评估连贯性时，可以指示评估者阅读原文和摘要，比较要点，核查逻辑顺序，再打分。
\item \textbf{评分函数。}大语言模型结合提示、思维链、原文与目标文本执行评估。评分函数根据模型回答评估提示时返回词元的概率给出分数。
\begin{equation}
\operatorname{score}=\sum_{i=1}^{n}p(s_i)\times s_i,
\label{eq:ch08-geval}
\end{equation}
其中，$s_i$ 为提示中预先规定的各个可能分值，$p(s_i)$ 为大语言模型给该分值分配的概率。
\end{enumerate}

\subsection{Pass@k}
\refcite{ch08-ref-037}{Kulal 等人（2019）}最早提出 pass@$k$，可用来评估语言模型生成代码的功能正确性。对于每个编程问题，它用单元测试衡量多次尝试中至少生成一份正确代码的可能性。实施步骤如下。

\begin{enumerate}
\item 为每道题生成 $n$ 份代码样例，$n\geq k$。例如，$n=200$ 且 $k\leq100$。
\item 统计成功通过单元测试的样例数，记作 $c$，其中 $c\leq n$。
\item 用下式计算 pass@$k$ 的无偏估计量。
\begin{equation}
\mathrm{pass}@k=1-\frac{\binom{n-c}{k}}{\binom{n}{k}}.
\label{eq:ch08-pass}
\end{equation}
\end{enumerate}

计算结果表示，从 $n$ 份生成样例中选 $k$ 份，至少一份通过单元测试的概率，可用来衡量模型解决编程任务的能力。

% 原书第 8.4 节。
\section{大语言模型的基准数据集}
\label{sec:ch08-benchmarks}
基准通过一套标准化指标，使不同大语言模型之间的比较更公平。借助基准，可以找出在特定情境中表现最好的模型，跟踪同一个模型随时间取得的改进，也可以比较不同模型的长短。基准数据集可以由多项任务组成，也可以只包含一项任务。每项任务都有自己的评估指标，用于考察模型的特定语言能力。本节介绍若干重要数据集及其用途。


\medskip\noindent\textbf{多任务或通用能力}\par\nobreak
\textbf{基准}包括 MMLU（\refcite{ch08-ref-029}{Hendrycks et al., 2020}）、SuperGLUE（\refcite{ch08-ref-096}{Wang et al., 2019}）、BIG-bench（\refcite{ch08-ref-091}{Srivastava et al., 2022}）、GLUE（\refcite{ch08-ref-095}{Wang et al., 2018}）、BBH（\refcite{ch08-ref-091}{Srivastava et al., 2022}）、Blended Skill Talk（\refcite{ch08-ref-090}{Smith et al., 2020}）和 HELM（\refcite{ch08-ref-049}{Liang et al., 2022}）。
\par\textbf{用途}是评估语言模型在多种任务上的表现，综合考察通用语言理解、推理、生成等能力。

\medskip\noindent\textbf{语言理解}\par\nobreak
\textbf{基准}包括 CoQA（\refcite{ch08-ref-078}{Reddy et al., 2019}）、WiC（\refcite{ch08-ref-071}{Pilehvar and Camacho-Collados, 2018}）、Wikitext103（\refcite{ch08-ref-056}{Merity et al., 2016}）、PG19（\refcite{ch08-ref-073}{Rae et al., 2019}）、QQP（\refcite{ch08-ref-043}{Le et al., 2021}）、CB（De \refcite{ch08-ref-022}{Marneffe et al., 2019}）、CNSS（\refcite{ch08-ref-051}{Liu et al., 2018}）、CKBQA（\refcite{ch08-ref-046}{Li et al., 2016}）、AQuA（\refcite{ch08-ref-050}{Ling et al., 2017}）、OntoNotes（\refcite{ch08-ref-099}{Weischedel et al., 2011}）、HeadQA（\refcite{ch08-ref-094}{Vilares and Gómez-Rodríguez, 2019}）和 Twitter Dataset（\refcite{ch08-ref-009}{Blodgett et al., 2016}）。
\par\textbf{用途}是考察语言理解的不同方面，包括问答、语境中的词义消歧和情感分析。

\medskip\noindent\textbf{故事填空与句子补全}\par\nobreak
\textbf{基准}包括 StoryCloze（\refcite{ch08-ref-060}{Mostafazadeh et al., 2016}）、LAMBADA（\refcite{ch08-ref-066}{Paperno et al., 2016}）、AdGen（\refcite{ch08-ref-088}{Shao et al., 2019}）和 E2E（\refcite{ch08-ref-063}{Novikova et al., 2017}）。
\par\textbf{用途}是测试模型补全故事和句子的能力。完成这类任务需要理解叙事语境，进行常识推理，并生成连贯文本。

\medskip\noindent\textbf{物理知识与世界知识}\par\nobreak
\textbf{基准}包括 PIQA（\refcite{ch08-ref-008}{Bisk et al., 2020}）、TriviaQA（\refcite{ch08-ref-034}{Joshi et al., 2017}）、ARC、ARC-Easy 和 ARC-Challenge（\refcite{ch08-ref-018}{Clark et al., 2018}）、PROST（\refcite{ch08-ref-005}{Aroca-Ouellette et al., 2021}）、OpenBookQA（\refcite{ch08-ref-058}{Mihaylov et al., 2018}）和 WebNLG（\refcite{ch08-ref-026}{Ferreira et al., 2020}）。
\par\textbf{用途}是考察模型对物理概念和一般世界知识的理解，常采用问答形式。

\medskip\noindent\textbf{语境中的语言理解}\par\nobreak
\textbf{基准}包括 RACE、RACE-Middle 和 RACE-High（\refcite{ch08-ref-041}{Lai et al., 2017}）、QuAC（\refcite{ch08-ref-015}{Choi et al., 2018}）、StrategyQA（\refcite{ch08-ref-028}{Geva et al., 2021}）和 Quiz Bowl（\refcite{ch08-ref-011}{Boyd-Graber et al., 2012}）。
\par\textbf{用途}是评估模型在语境中理解和解释语言的能力。这是聊天机器人和内容分析等应用所需的能力。

\medskip\noindent\textbf{常识推理}\par\nobreak
\textbf{基准}包括 WinoGrande（\refcite{ch08-ref-082}{Sakaguchi et al., 2021}）、HellaSwag（\refcite{ch08-ref-109}{Zellers et al., 2019}）、COPA（\refcite{ch08-ref-080}{Roemmele et al., 2011}）、WSC（\refcite{ch08-ref-044}{Levesque et al., 2012}）、CSQA（\refcite{ch08-ref-092}{Talmor et al., 2018}）、SIQA（\refcite{ch08-ref-083}{Sap et al., 2019}）和 ReCoRD（\refcite{ch08-ref-110}{Zhang et al., 2018}）。
\par\textbf{用途}是考察模型能否在复杂的自然语言任务中运用常识推理、因果理解和现实世界知识。

\medskip\noindent\textbf{阅读理解}\par\nobreak
\textbf{基准}包括 SQuAD（\refcite{ch08-ref-074}{Rajpurkar et al., 2016}）、BoolQ（\refcite{ch08-ref-016}{Clark et al., 2019}）、SQuADv2（\refcite{ch08-ref-075}{Rajpurkar et al., 2018}）、DROP（\refcite{ch08-ref-025}{Dua et al., 2019}）、RTE（\refcite{ch08-ref-021}{Dagan et al., 2005}）、WebQA（\refcite{ch08-ref-013}{Chang et al., 2022}）、MultiRC（\refcite{ch08-ref-035}{Khashabi et al., 2018}）、Natural Questions（\refcite{ch08-ref-039}{Kwiatkowski et al., 2019}）、SciQ（\refcite{ch08-ref-100}{Welbl et al., 2017}）和 QA4MRE（\refcite{ch08-ref-070}{Peñas et al., 2013}）。
\par\textbf{用途}是测试模型分析和理解文段，并根据文段回答问题的能力。

\medskip\noindent\textbf{数学推理}\par\nobreak
\textbf{基准}包括 MATH（\refcite{ch08-ref-030}{Hendrycks et al., 2021}）、Math23k（\refcite{ch08-ref-097}{Wang et al., 2017}）、GSM8K（\refcite{ch08-ref-019}{Cobbe et al., 2021}）、MathQA（\refcite{ch08-ref-006}{Austin et al., 2021}）、MGSM（\refcite{ch08-ref-089}{Shi et al., 2022}）、MultiArith（\refcite{ch08-ref-081}{Roy and Roth, 2016}）、ASDiv（\refcite{ch08-ref-057}{Miao et al., 2021}）、MAWPS（\refcite{ch08-ref-036}{Koncel-Kedziorski et al., 2016}）和 SVAMP（\refcite{ch08-ref-068}{Patel et al., 2021}）。
\par\textbf{用途}是评估模型解答数学问题的能力，范围从基础算术到更复杂的代数和几何问题。

\medskip\noindent\textbf{问题求解}\par\nobreak
\textbf{基准}包括 HumanEval（\refcite{ch08-ref-014}{Chen et al., 2021}）、DS-1000（\refcite{ch08-ref-042}{Lai et al., 2023}）、MBPP（\refcite{ch08-ref-006}{Austin et al., 2021}）、APPS（\refcite{ch08-ref-030}{Hendrycks et al., 2021}）和 CodeContests（\refcite{ch08-ref-048}{Li et al., 2022}）。
\par\textbf{用途}是测试模型运用逻辑和推理解决各种问题的能力，其中包括编程挑战。

\medskip\noindent\textbf{自然语言推断与逻辑推理}\par\nobreak
\textbf{基准}包括 ANLI（\refcite{ch08-ref-062}{Nie et al., 2019}）、MNLI-m 和 MNLI-mm（\refcite{ch08-ref-102}{Williams et al., 2017}）、QNLI（\refcite{ch08-ref-074}{Rajpurkar et al., 2016}）、WNLI（\refcite{ch08-ref-044}{Levesque et al., 2012}）、ANLI R1、ANLI R2 和 ANLI R3（\refcite{ch08-ref-062}{Nie et al., 2019}）、HANS（\refcite{ch08-ref-055}{McCoy et al., 2019}）、LogiQA（\refcite{ch08-ref-052}{Liu et al., 2020}）和 StrategyQA（\refcite{ch08-ref-028}{Geva et al., 2021}）。
\par\textbf{用途}是考察模型根据给定文本作出推断的能力。这是理解自然语言和用自然语言推理的重要组成部分。

\medskip\noindent\textbf{跨语言理解}\par\nobreak
\textbf{基准}包括 MLQA（\refcite{ch08-ref-045}{Lewis et al., 2019}）、XNLI（\refcite{ch08-ref-020}{Conneau et al., 2018}）、PAWS-X（\refcite{ch08-ref-104}{Yang et al., 2019}）、XSum（\refcite{ch08-ref-061}{Narayan et al., 2018}）、XCOPA（\refcite{ch08-ref-072}{Ponti et al., 2020}）、XWinograd（\refcite{ch08-ref-093}{Tikhonov and Ryabinin, 2021}）、TyDiQAGoldP（\refcite{ch08-ref-017}{Clark et al., 2020}）和 MLSum（\refcite{ch08-ref-086}{Scialom et al., 2020}）。
\par\textbf{用途}是评估模型理解和处理不同语言语境的能力，对多语言环境中的应用很重要。

\medskip\noindent\textbf{语言翻译}\par\nobreak
\textbf{基准}包括 WMT（\refcite{ch08-ref-010}{Bojar et al., 2016}）、WMT20 和 WMT20-enzh（\refcite{ch08-ref-054}{Loïc et al., 2020}）。
\par\textbf{用途}是评估模型在不同语言之间翻译文本的熟练程度。翻译是自然语言处理的一项基本任务。

\medskip\noindent\textbf{对话}\par\nobreak
\textbf{基准}包括 Wizard of Wikipedia（\refcite{ch08-ref-023}{Dinan et al., 2018}）、Empathetic Dialogues（\refcite{ch08-ref-077}{Rashkin et al., 2018}）、DPC 生成的对话（\refcite{ch08-ref-031}{Hoffmann et al., 2022}）和 ConvAI（\refcite{ch08-ref-024}{Dinan et al., 2020}）。
\par\textbf{用途}是评估模型进行连贯、符合语境的对话的能力，这对开发聊天机器人十分重要。

% 原书第 8.5 节。
\section{大语言模型选型}
\label{sec:ch08-selection}
为应用选出最合适的大语言模型，可以说是其中最重要的一项决策。语言模型及其建模方法的能力不断提高，OpenAI 在 2022 年 11 月发布 ChatGPT 是一个显著节点。这是撰写本书的主要动因，也使采用这些模型的领域涌现出大量创新。不过，选用哪一种模型时，模型能力或分析质量只是一项考量。

\begin{keypoint}
\textbf{为模型选型和定制方式绘制决策流程图}

以下各节讨论大语言模型的各种属性及其与应用开发情境的关系。读者可以自行绘制决策流程图，把模型属性对应到应用领域。这样做有助于解决应用所针对的问题，使方案日后仍可维护，并在提高效率和效果的同时，减少采用大语言模型带来的风险。请仔细考虑应用的功能需求和非功能需求，将其与下文讨论的机会和难题逐项联系起来，据此建立决策过程。
\end{keypoint}

模型选择发生在项目早期，会影响许多后续选项，因此还应考虑许多其他标准和属性。图~\ref{fig:ch08-selection}改绘自 \refcite{ch08-ref-047}{Li et al. (2023)}，以一个领域的选型和开发过程为例，说明如何在开源和闭源模型之间作出选择，并根据工具、数据和预算等条件选择定制路径。定制方法自上而下按成本由低到高排列，便于在考虑开销时依次尝试。

\begin{figure}[!htbp]
\centering
\begin{tikzpicture}[
  x=1cm,y=1cm,font=\sffamily\footnotesize,
  condition/.style={ellipse,draw=cyan!50!black,fill=cyan!13,align=center,
    text width=2.33cm,minimum width=3.3cm,minimum height=.78cm,inner sep=1pt},
  choice/.style={rectangle,rounded corners=2pt,draw=orange!65!black,fill=orange!15,
    align=center,text width=2.57cm,minimum height=.83cm,inner sep=3pt},
  level/.style={align=center,text width=1.16cm,font=\sffamily\scriptsize,inner sep=0pt},
  yes/.style={-{Stealth[length=1.7mm]},draw=green!50!black},
  no/.style={-{Stealth[length=1.7mm]},draw=red!70!black},
  next/.style={-{Stealth[length=1.7mm]},draw=black!80},
  route/.style={font=\sffamily\scriptsize,fill=white,inner sep=1pt}]
\node[condition] (private) at (0,0) {数据需要\\保密吗？};
\node[condition] (examples) at (0,-1.78) {能提供若干\\任务示例吗？};
\node[condition] (tools) at (0,-3.56) {有工具可辅助\\注入知识吗？};
\node[condition] (domain) at (0,-5.34) {有领域数据\\可用于微调吗？};
\node[condition] (budget) at (0,-7.28) {预算逾 100 万\\美元且有数十亿\\训练词元吗？};
\node[choice] (own) at (-3.40,-.79) {1A 自行部署\\开源大语言模型};
\node[choice] (api) at (3.40,-.79) {1B 使用第三方\\大语言模型（如 GPT-4）};
\node[choice] (few) at (-3.40,-2.57) {2A 对选定模型\\采用少样本方法};
\node[choice] (rag) at (-3.40,-4.35) {3A 用工具增强\\所选大语言模型};
\node[choice] (tune) at (-3.40,-6.13) {3B 通过微调\\向模型注入知识};
\node[choice] (train) at (-3.40,-8.25) {4A 从头训练\\自己的大语言模型};
\node[draw,rounded corners,fill=black,text=white,font=\sffamily\scriptsize] (end) at (3.40,-8.25) {无可行方案};
\draw[yes] (private.west) to[out=180,in=90] node[route,above,pos=.52]{是} (own.north);
\draw[no] (private.east) to[out=0,in=90] node[route,above,pos=.52]{否} (api.north);
\draw[no] (private) -- node[route,right]{否} (examples);
\draw[next] (own.south) to[out=-90,in=180] (examples.north west);
\draw[next] (api.south) to[out=-90,in=0] (examples.east);
\foreach \question/\option/\following in {examples/few/tools,tools/rag/domain,domain/tune/budget}{
  \draw[yes] (\question.south west) to[out=180,in=90] node[route,above,pos=.52]{是} (\option.north);
  \draw[no] (\question) -- node[route,right]{否} (\following);
  \draw[next] (\option.south) to[out=-90,in=180] (\following.north west);
}
\draw[yes] (budget.south west) to[out=180,in=90] node[route,above,pos=.52]{是} (train.north);
\draw[no] (budget.east) to[out=0,in=90] node[route,above,pos=.52]{否} (end.north);
\draw[next] (train.east) -- (end.west);
\node[level] at (-5.42,-.79) {第 1 级\\零样本};
\node[level] at (-5.42,-2.57) {第 2 级\\少样本};
\node[level] at (-5.42,-5.24) {第 3 级\\RAG 或\\微调};
\node[level] at (-5.42,-8.25) {第 4 级\\从头训练};
\draw[next] (2.12,-2.82) -- (2.82,-2.82) node[right,font=\sffamily\scriptsize]{效果不满意};
\node[font=\sffamily\scriptsize,draw=orange!65!black,fill=orange!15,rounded corners=2pt]
  at (3.29,-3.47) {决策框};
\node[font=\sffamily\scriptsize,ellipse,draw=cyan!50!black,fill=cyan!13]
  at (3.29,-4.17) {方案框};
\draw[-{Stealth[length=3mm]},line width=1.25pt,draw=orange!75!black]
  (3.29,-4.77) -- node[right,font=\sffamily\scriptsize]{成本上升} (3.29,-6.85);
\end{tikzpicture}
\caption{金融服务领域中开源与闭源大语言模型的选型及定制流程。}
\label{fig:ch08-selection}
\end{figure}


\begin{keypoint}
为项目选型时，总拥有成本（total cost of ownership，TCO）也是一种有用的分析框架。它按项目的模型、使用场景等细节估算各项成本，再加总以便比较方案。部分项目如下。
\begin{itemize}
\item \textbf{每词元成本}，反映每次查询的处理和生成开销。
\item \textbf{人力成本}，估算构建和部署大语言模型服务所需的人力开销。
\item \textbf{整体搭建成本}，估算部署、运行和管理大语言模型服务的总开销。
\end{itemize}
希望进一步了解这些因素及计算过程的开发人员，可以从 Hugging Face 的资料入手，其中包括一个交互式 TCO 计算器。读者也可以查阅类似资料，了解这一框架怎样辅助应用开发中的决策。
\end{keypoint}


\transnote{所据英文版标出了注 3 和注 4 的序号，页面上未印出这两条注释的链接。}

各类大语言模型各有利弊，哪些利弊重要很大程度上取决于应用情境。例如，一个应用主要做文本摘要，选用号称在实体抽取基准上达到当时最佳水平的模型，意义不大。本节介绍选型的主要标准，帮助开发人员判断模型属性是否切合其应用，从而作出选择。

\subsection{开源与闭源}
开源还是闭源，是开发人员为应用选择大语言模型时首先考虑的问题之一。两者的主要取舍通常在使用灵活性、便利性和成本三个方面，还应考虑表~\ref{tab:ch08-open-closed}所列的其他因素。例如，开发团队起初可能看中开源模型的使用成本低、灵活性高，但按表中各项标准逐一评估后，有可能发现采用开源模型的应用，其总拥有成本远高于采用闭源模型的应用。因此，开发人员应尽量全面地评估模型选择。以下各项分别讨论开源和闭源模型的有关倾向。

{\small\setlength{\tabcolsep}{4pt}\renewcommand{\arraystretch}{1.12}
\begin{longtable}{@{}>{\raggedright\arraybackslash}p{15mm} >{\raggedright\arraybackslash}p{47mm} >{\raggedright\arraybackslash}p{47mm}@{}}
\caption{应用开发中选择开源或闭源大语言模型时应考虑的方面。}\label{tab:ch08-open-closed}\\
\toprule
方面 & 闭源 & 开源\\
\midrule
\endfirsthead
\multicolumn{3}{@{}l}{表 \thetable\ （续）}\\
\toprule
方面 & 闭源 & 开源\\
\midrule
\endhead
\midrule
\multicolumn{3}{r@{}}{下页续}\\
\endfoot
\bottomrule
\endlastfoot
可获取性 & 终端用户往往无法获取模型架构、训练数据、学习目标及源代码的详细信息。 & 终端用户可依宽松许可获取模型架构、训练数据、学习目标和源代码的详细信息。\\
透明度 & 终端用户通常通过预先规定的 API 使用模型，上述各项的透明度有限，影响在不同情境中的使用，也影响应用的调试、维护和支持。 & 终端用户通常可获得模型权重。开发人员负责托管模型并提供服务。透明度高，支持多种定制和灵活的使用场景，相应支持也由开发人员承担。\\
成本 & 使用成本通常高于开源模型，但便利性和使用支持等因素往往可以抵消较高的价格。 & 使用成本可能远低于闭源模型，部分原因是终端用户可以定制并优化计算资源消耗等。较低的使用成本应与开发和支持开销一并权衡，后两项均由终端用户承担。\\
定制 & 仅限于规定的方法。这些方法的成本往往高于开源方案。 & 可灵活采用适应性预训练、全参数微调、参数高效微调或基于提示的学习。\\
使用方式 & 通常限于若干预先设定的方式，例如图形用户界面或 API。 & 使用方式灵活，可由终端用户自行设计和开发。\\
协作 & 仅限于已集成的协作工具。 & 协作方式灵活，可采用终端用户所选的任何协作工具或框架。\\
安全 & 由模型所有者全面管理。这对没有专门安全资源的小型开发团队很有价值。然而，任何安全方案都无法做到百分之百有效，闭源条件也可能使模型或服务中的漏洞难以发现。 & 终端用户必须自行开发和管理安全措施。这样可获得最大的灵活性，但对小型开发团队而言，承担这项责任可能复杂且昂贵。\\
隐私 & 输入模型的任何数据均受终端用户协议约束。提供方可能获准用这些数据改进未来版本，因此还要考虑知识产权流失。 & 终端用户根据需求或偏好决定隐私保护的程度。\\
支持 & 终端用户可获得模型所有者的支持，并负责为自己应用中封装的功能提供支持。 & 模型往往不附带终端用户支持。活跃的开源社区有时可通过 GitHub 或 Hugging Face 项目为用户解决问题提供帮助。\\
更新与维护 & 模型所有者安排并发布更新。更新过程可能不透明，并受模型所有者的商业规划支配，未必符合所有终端用户的意愿。 & 终端用户负责全部模型更新与维护。这样可获得最大的控制权，但对较小的开发团队而言，可能需要承担较高的开销。\\
\end{longtable}
}

\subsection{分析质量}
在模型选型的各项考量中，权重最高的通常是模型完成所要解决的使用场景相关任务的质量。较大的模型往往具有较强的分析表现，因而计算成本和财务开销是主要的取舍。汇总的排行榜可作为衡量分析质量的参考。排行榜用一系列标准化基准评估大语言模型，便于直接比较和筛选。解读结果时应仔细考虑使用的评估方法（\refcite{ch08-ref-003}{Alzahrani et al., 2024}）。

即便如此，排行榜也可以大致显示模型之间的相对表现，有助于迅速缩小候选范围，再进一步评估各模型是否适合项目。原书认为第~\ref{sec:ch08-categories}~节所述的 Hugging Face Open LLM Leaderboard 在当时维护得最好，但只覆盖开源模型。也有同时覆盖开源和闭源模型的排行榜，但这些项目能否长期维护并不明确，例如 \url{https://llm-leaderboard.streamlit.app/}。希望同时评估两类模型时，可能需要检索合适的资料。

把候选模型缩小到可管理的范围后，可以用更有针对性的任务进一步评估分析表现。大语言模型的任务表现对所用数据敏感。换用更接近应用领域的数据分布，模型的实际表现可能低于排行榜上的基准结果。第二轮贴近使用场景的评估可进一步缩小范围，找出在该评估中表现最好的模型，或在项目容许继续开展提示工程、预训练或微调时有改进潜力的模型（\refcite{ch08-ref-108}{Yuan et al., 2023}）。

\subsection{推理延迟}
\label{sec:ch08-latency}
大语言模型的推理延迟是选型的重要因素，可以理解为从处理单个请求到把对应响应返回给用户所需的总时间，见图~\ref{fig:ch08-latency}。把用户的输入提示传到模型 API 的延迟（图中第~1~步），以及把模型响应从 API 送回用户的延迟（第~4~步），大部分可通过网络优化来处理。除此之外，对总体延迟影响最大的有两个推理阶段，即模型网络处理输入提示中词元的时间（第~2~步），以及抽样和组织响应词元的时间（第~3~步）。前者称为预填充（prefill），后者称为解码（decode）（\refcite{ch08-ref-002}{Agrawal et al., 2024}）。

\sourcefigure{fig08_03.pdf}{大语言模型的延迟由四部分组成。（1）用户提示到达模型推理 API 的时间。（2）模型处理用户提示的时间。（3）模型抽取相关词元并组成响应的时间。（4）模型响应从 API 传给用户的时间。}{fig:ch08-latency}

\urldef\latencytableurl\url{https://huyenchip.com/2023/04/11/llm-engineering.html#cost_and_latency}
由于采用 Transformer 架构，预填充阶段可并行处理提示词元，因此延迟相对较高（与解码步骤相比），计算资源利用率也高。相较而言，解码是顺序过程，因为在生成下一个输出词元之前，必须先生成所有之前的词元。因此，单个输出词元的延迟相对较低，顺序生成也使计算资源利用率较低。这意味着与输出长度相比，提示中的输入词元数对推理延迟的影响不大。例如，表~\ref{tab:ch08-latency}\footnote{表中数据转载自 \latencytableurl。}给出了调整输入和输出词元长度对 OpenAI \texttt{gpt-3.5-turbo} 模型响应延迟的影响。把输入词元从 51 个增至 232 个，而输出词元维持为 1 个，延迟变化很小。输入长度大致相同时，若把输出从 1 个词元增至 26 个，延迟接近原来的 3 倍。这显示输入和输出长度对推理延迟的影响并不均衡。



\begin{table}[htbp]
\centering
\caption[输入和输出长度对推理延迟的影响]{输入和输出长度对推理延迟的影响。数据来自 OpenAI 的 \texttt{gpt-3.5-turbo} 模型。OpenAI 如何调度和路由用户查询，用户无从得知，因此部分波动来自 API 延迟。即便绝对延迟发生变化，输入和输出长度设置之间的关系仍保持一致。p50 延迟以秒计，表示在 $n=20$ 次请求中，有 50\% 的请求响应时间不超过该值。}
\label{tab:ch08-latency}
\begin{tabular}{rrr}
\toprule
输入词元数 & 输出词元数 & p50 延迟（秒）\\
\midrule
51 & 1 & 0.58\\
232 & 1 & 0.53\\
228 & 26 & 1.43\\
\bottomrule
\end{tabular}
\end{table}

考虑到这种不均衡，大语言模型有哪些属性会影响推理延迟？最直观的是模型规模。经验上，参数越多，延迟越高。参数更多的模型需要更多计算来处理输入并生成输出。模型架构也很重要。层数、各层的复杂程度、Transformer 块所用的注意力机制，以及 Transformer 块在网络中的数量和位置，都会影响推理延迟。

参数以何种数值精度存储，也是影响大语言模型推理延迟的重要因素。第~\ref{ch:adaptation}~章第~\ref{sec:ch04-quantization}~节已详细讨论量化。在比较开源和闭源模型时，两者在定制方面的差别值得注意。闭源模型定制的限制较多，终端用户只能采用模型所有者支持的量化方式。开源模型的终端用户通常可以自行测试并采用最适合使用场景的量化方法。量化有可能显著降低推理延迟，以及运行或托管模型所需的内存和存储开销，因此应认真考虑闭源模型在定制上的限制。如果预计请求与响应的次数很少，这一点可能不太重要。然而，在请求与响应次数较多的情况下，闭源大语言模型可能成为应用的瓶颈。例如，按原书当时的描述，OpenAI 的 API 通常按端点和模型分别设定速率限制。

\subsection{成本}
应用开发使用大语言模型时，很多环节都要付出成本。人们往往只考虑推理的每词元成本，它确实是最重要的成本之一。不过，大量研究和商业投入不断压低这项成本，其数值也随之改变。例如，图~\ref{fig:ch08-costs}给出了当时的情况。\footnote{转载自 \url{https://medium.com/@ja_adimi/comparison-cost-analysis-should-we-invest-in-open-source-or-closed-source-llms-bfd646ae1f74}。}围绕开源与闭源模型的讨论通常认为，闭源模型每个词元的成本显著高于开源模型。不过，这很可能只适用于能力最强的闭源模型版本。原书以 OpenAI 当时的定价为例，较旧型号的推理开销通常只是最新型号的一小部分。

\sourcefigure{fig08_04.pdf}{生成 100 万个词元的美元成本。图中比较 OpenAI 的两个闭源大语言模型 GPT-4 和 GPT-3.5 Turbo，以及在 A100 和 A10G 两种 GPU 上运行的开源模型 Llama-2-7B。两种配置的 Llama-2 生成成本均低于 OpenAI 当时能力最强的模型 GPT-4 的十分之一，但 GPT-3.5 Turbo 比这两种配置都便宜。选择项目所用模型时，应考虑闭源模型推理成本下降的趋势。费用为 2023 年 8 月分析时的数值。}{fig:ch08-costs}

采用总拥有成本框架时，除了每词元生成成本，开发人员也应估算应用开发生命周期各环节的开销。例如，采用开源模型的推理成本可能更低，分析质量也可能较低。微调或投入提示工程也许能缩小这种差距，但优化本身要付出成本。微调或测试都需要数据，收集、标注和整理数据可能费力且复杂。此外，模型经微调后还需要长期维持其表现，这意味着持续评估并在必要时改进模型的工作也应计入成本。定制和维护大语言模型是复杂的技术工作，因而相较于采用闭源模型，项目需要承担更高的专业技术与人才成本。选择开源模型时，托管和计算资源管理也是直接开销，会增加整体复杂度和成本。

\subsection{适应性与维护}
开源大语言模型通常更易适配，因为终端用户往往可以直接获得模型权重、训练数据和源代码。这使用户可以采用第~3~章和第~4~章介绍的一项或多项技术来适配或定制模型，进而控制应用的行为和表现。灵活性的代价是开发生命周期需要消耗更多资源。

闭源模型适应性较低，相应地，其开发周期消耗的资源也少得多。如果项目利用大语言模型执行常见任务，专有模型很可能已经具备较好的开箱能力，无需再适配或定制。类似地，采用 $n$ 样本上下文学习等高级提示工程技术，还可以进一步改善结果。这些情况下可以考虑闭源或专有模型。持续维护仍然是决策因素。闭源模型的维护通常不透明，会按照模型所有者的规划或维护日程进行。假如变更可以不事先通知终端用户，开发人员就应了解所选模型悄然升级对应用表现的影响。升级后的模型会降低用户体验吗？如果处置不当，是否会引发伦理或安全风险？

通过在开发生命周期中引入持续监控和评估，可以减少相当一部分这类风险。不过，大语言模型的规模和复杂性使得事先预测所有失效情形成为不可能。因此，选用何种大语言模型，取决于具体使用场景。在应用出现错误的代价很高的环境中（例如受监管行业），可能无法任由用户的结果碰运气，或由 OpenAI、Anthropic 等第三方决定。此时，应用提供方只能选择自己能够全面负责这些风险的模型。

\subsection{数据安全与许可}
应用常借助来自用户或其他来源的敏感数据提供结果。如果把敏感数据作为上下文写进提示并发送给大语言模型，很难彻底阻止模型响应包含敏感上下文，因此会出现数据安全或隐私风险。许多方案采用先生成再过滤的方法，在向用户提供响应之前删除其中的敏感数据。也可采用第~5~章综述的对齐方法，减少用户收到敏感或不希望出现的信息的风险。

对于闭源模型或由第三方托管的模型，数据经网络传送到模型 API 会带来额外的安全问题。如果模型为专有模型，还须考虑所有者如何留存和使用数据。ChatGPT 刚向公众开放时，OpenAI 如何利用服务中突然涌入的用户查询和交互，是一个热门话题。许多用户担心 OpenAI 会把他们的提示用于进一步改进 ChatGPT，其中可能包含私密或敏感信息，进而可能暴露给其他用户。由于 ChatGPT 的流行程度增长很快，这个问题当时并不明朗。关于通过越狱从大语言模型提取“训练数据”的各种报告，使医疗和金融服务等对数据隐私及安全要求很高的行业中，许多潜在用户心存疑虑（\refcite{ch08-ref-107}{Yu et al., 2024}）。

规划应用设计及维护工作时，开发人员应了解许可方面的两个问题。其一与传统软件和工具的许可相似，即大语言模型发布时所用的许可往往规定了应用范围。例如，某些许可只准用于研究或实验，另一些可能使应用的营销或推广比预期复杂。本节不详细讨论不同的许可类型及其对应用开发人员的影响，感兴趣的读者可查阅有关资料。\footnote{\url{https://github.com/eugeneyan/open-llms}}

其二是预训练数据的许可。越来越多的材料表明，至少部分模型训练数据可能受版权限制，把这些材料纳入训练集可能构成侵权。备受关注的案例之一，是美国法院审理的 \textit{The New York Times} 与 OpenAI 之间的案件。在这起案件中，《纽约时报》主张 OpenAI 未经授权将该媒体的内容纳入大语言模型，侵犯了其权利（\refcite{ch08-ref-087}{S.D.N.Y., 2023}）。在此类案件的审理过程中，应用开发人员需预判裁决对应用维护的影响。如果法院支持著作权主张方，模型所有者可能须撤下使用受版权保护的材料训练的模型，并发布更新版本。应用开发人员因而可能不得不更换应用中的模型，并补做测试以确认表现是否一致。

% 原书第 8.6 节。
\section{应用开发工具}
\label{sec:ch08-tooling}
大语言模型领域迅速创新，配套工具生态也随之扩张。许多工具专用于微调模型、优化提示等特定开发阶段，另一些功能更丰富，可以在多个阶段发挥作用。对不熟悉其演变的人来说，在这个生态中选择工具并不容易。表~\ref{tab:ch08-tools}列举了部分常用工具，可以看出创新的范围。\footnote{据 \url{https://github.com/langgenius/dify} 修改。}

\begin{table}[htbp]
\centering\small
\caption{为构建和部署大语言模型应用提供支持的部分工具。}
\label{tab:ch08-tools}
\begin{tabularx}{\linewidth}{@{}p{36mm}>{\raggedright\arraybackslash}X@{}}
\toprule
在大语言模型应用中的作用 & 常用工具\\
\midrule
嵌入与索引 & Pinecone、Weaviate、Milvus、Chroma\\
数据标注 & Scale、Snorkel、Label Studio\\
训练、开发与评估 & Hugging Face、Lightning AI、HumanLoop\\
实验跟踪 & Weights \& Biases、MLflow、Comet\\
隐私、安全与合规 & Lakera、Skyflow、Nomic\\
监控与可观测性 & TruLens、WhyLabs、Arize\\
托管与部署 & Lambda、Together.ai、Groq、Predibase、Anyscale\\
提示串联与集成 & LlamaIndex、LangChain、DSPy\\
\bottomrule
\end{tabularx}
\end{table}

本节按功能介绍这一生态。首先介绍若干用于连接大语言模型应用各组件的重要工具。这些工具通常采用链式结构。与一个或多个模型的交互，以及应用设计中的其他组件，都可以模块化并依次串联起来，以支持丰富的工作流和用户交互。随后介绍大语言模型定制工具，具体考察预训练、适应性预训练和微调所用的库，指出其中较常用的工具和具备独特能力的工具。接下来讨论提示工程及其可选工具，再介绍开发人员可用的若干向量数据库，以及便于集成它们的工具。最后讨论应用开发中的评估与监控。

\subsection{大语言模型应用框架}
\label{sec:ch08-frameworks}
大语言模型应用框架用来连接应用中数量往往很多的组件。这些工具通常对应用开发方式有一定约束，因此选择的框架应与应用所需的开发模式相符。有些框架的功能比另一些丰富。例如，表~\ref{tab:ch08-frameworks}从八项特征比较当时的三个框架，其中只有 Dify.ai 支持集成单点登录（single sign-on，SSO）等企业功能。如果应用有这项要求，选择 Dify.ai 可能比在采用 LangChain 开发的应用中自行构建 SSO 更合适。

\begin{table}[htbp]
\centering\small
\caption{三个常用大语言模型应用开发框架的功能比较。}
\label{tab:ch08-frameworks}
\setlength{\tabcolsep}{4pt}
\begin{tabularx}{\linewidth}{@{}p{27mm}*{3}{>{\centering\arraybackslash}X}@{}}
\toprule
框架特征 & Dify.ai & LangChain & Flowise\\
\midrule
编程方式 & API 与面向应用 & Python 代码 & 面向应用\\
支持的模型 & 种类丰富 & 种类丰富 & 种类丰富\\
RAG 引擎 & 支持 & 支持 & 支持\\
智能体 & 支持 & 支持 & 支持\\
工作流 & 支持 & — & —\\
可观测性 & 支持 & 支持 & —\\
企业功能 & 支持 & — & —\\
本地部署 & 支持 & 支持 & 支持\\
\bottomrule
\end{tabularx}
\end{table}

较早出现且较流行的框架之一是 LangChain。\footnote{\url{https://github.com/langchain-ai/langchain}}这一开源项目旨在帮助开发人员更快、更可靠地把想法变成生产应用。它围绕 Python 构建，Python 也是与 Hugging Face 等框架交互最成熟的语言，因此适合用这门语言开发的应用。LangChain 的主要优点是控制范围广、适应性强。应用与大语言模型及用户输入的交互较为复杂时，这类框架通常最有价值。LlamaIndex\footnote{\url{https://www.llamaindex.ai/}}是另一种应用开发框架，旨在让开发人员通过较高层的接口使用较复杂的功能。如第~7~章所示，LlamaIndex 提供多种 RAG 应用模式与组件，适合采用这种范式的应用。在通用大语言模型应用开发中，它也是 LangChain 之外的合适选择。

Flowise\footnote{\url{https://github.com/FlowiseAI/Flowise}}是另一个值得关注的大语言模型应用开发框架。它基于 JavaScript 构建，用户可以在图形界面中拖放组件，以可视方式定制大语言模型工作流。无代码的流程设计和开发对于快速制作原型，或偏好低代码、无代码开发的项目尤其有用。Flowise 与 JavaScript 生态结合紧密，可以快速集成到 Web 应用中，从而缩短重视上市时间的项目获得收益所需的时间。

上述框架大多围绕较高层次的顺序工作流组织，DSPy\footnote{\url{https://github.com/stanfordnlp/dspy}}采用了另一种方法。它提供较底层的模块，用户可自行组合，实现与 LangChain 预构建模块类似的功能。DSPy 的一个特点是编译器，它能自动优化多步流水线，提高流水线在各项任务上的表现。对于数据输入、控制流程、执行次序乃至所用大语言模型都可能发生变化的应用，这种自动优化很有用。DSPy 的功能并不依赖预构建的提示模板等组件，适合需要从底层控制应用与大语言模型交互的开发人员。

最后介绍 Dify.ai。如表~\ref{tab:ch08-frameworks}所示，该框架功能较完整，适合企业级应用开发。它也支持 LLMOps，明确采用多种开发模式，使开发过程和应用行为可检验、可靠且可复现。对企业应用而言，这些模式十分重要，因为故障或低效率通常会造成更高的财务成本，也可能带来更大的社会影响，例如金融服务把大语言模型用作交易智能体时。

了解辅助开发的平台后，以下各节继续介绍可在框架内使用的工具，用来完成提示工程、大语言模型定制和评估等较具体的工作。

\subsection{大语言模型定制}
如前所述，闭源与开源大语言模型的定制方式差别很大。闭源模型通常按照所有者规定的流程定制。开源模型的用户一般可采用任何技术上可行的方法。两者的技术复杂度也不同。对于没有深厚机器学习专业知识的应用开发团队，闭源模型的定制是可选方案。定制开源模型需要更多技术能力。

OpenAI 当时通过 API 提供基础模型微调，开发人员可借助其软件开发工具包（SDK）库与 API 交互。微调流程包括准备数据集，按 OpenAI 的 Chat Completions API 规定设置数据格式，并划分训练集和测试集，以便在微调中评估模型。随后通过 OpenAI API 上传数据，启动微调。在这一过程中，用户可试验不同的超参数、数据质量和格式，迭代改善微调结果。开发人员应仔细考虑迭代的成本，因为费用由所用训练词元数、待微调的 OpenAI 模型的基础词元单价和训练轮数等多项因素共同决定。

开源大语言模型的定制选择范围更广，复杂度也高得多。如果没有弄清其复杂度，试验中累积的额外算力开销可能使微调成本大幅增加。AWS Bedrock 或 Google Vertex AI 等云服务可为开发人员管理模型所需的计算资源。也有 AWS SageMaker JumpStart 等自主管理方案。采用后者时，开发人员大多可以自行选择微调或对齐工具包。PyTorch 等工具是这条定制路径的常用基础，Hugging Face Transformers 和 Hugging Face PEFT 等更高层工具可为其提供封装。围绕微调复杂性、计算和成本效率也涌现出越来越多的专用工具，例如 ggml 和 LLMZoo。有关定制的细节和使用示例，可重读第~4~章和第~5~章的教程。

\subsection{向量数据库}
语言模型的能力刚开始快速增强时，语义相似性很快成为其突出的用途之一。将两段文本分别输入模型，比较它们的嵌入，即可判断二者语义是否相似。通常可用两个嵌入之间的余弦距离。不过，如果要从大型知识库中找出最相似的文本，现场运行模型并计算数百万对向量的余弦距离会耗费很长时间。解决方法是预先计算每个文本块的嵌入，将它们存入向量数据库，以便高效检索。

Facebook 的 FAISS（\refcite{ch08-ref-033}{Johnson et al., 2017}）是大规模向量搜索早期的成功实例之一，也是一个包含多种索引技术的开源库。此后出现了许多向量数据库方案，例如 Pinecone、Milvus 和 Chroma。向量数据库针对向量搜索的存储和检索两个环节进行优化（\refcite{ch08-ref-085}{\refcite{ch08-ref-085}{Schwaber-Cohen, 2023}}）。Postgres 和 Cassandra 等既有数据库也开始支持向量存储。

向量数据库的使用场景广泛。最常见的场景之一是第~7~章的 RAG。RAG 应用的检索步骤往往要在向量数据库中找到回答用户所需的信息。另一个重要用途是问答。例如，客户对某种产品有一个常见问题，却在公司网站上找不到答案。如果其他用户问过类似问题，也许可以复用已有答案，无需再由人工查找。更一般地说，向量搜索可以补充传统关键词搜索。关键词搜索的结果较容易预判，且可具有较高召回率。向量搜索扩大了可能检出的文档范围，在一次检索中结合二者往往有效。

\subsection{提示工程}
第~3~章介绍过不同的提示形式。从更适合掩码语言模型的基本填空式、前缀式提示，到针对 GPT-4 或 Llama-2 等生成模型在连续空间优化过的提示，都属于提示工程的范畴。理论和方法可参阅前面的章节，本节重点介绍在应用开发项目中开发和维护提示的实用工具。

快速发展的除了大语言模型，还有直接使用模型并将其集成到应用中的配套工具。工具的质量往往到实际使用时才会显现，这给开发人员带来了困难。希望了解更多提示工程工具的读者可以查阅 \url{https://www.promptingguide.ai/tools} 和 \url{https://learnprompting.org/docs/tooling/tools}。这里选取若干常用工具，说明它们能提供哪些功能，以及可用哪些方式开发提示。

图~\ref{fig:ch08-prompt-workflow}在生产应用开发的整体背景下展示了提示工程的一种典型流程（\refcite{ch08-ref-007}{Benram, 2023}）。提示工程和改进通常依据若干评估标准开展，包括在基准或测试数据集上的分析表现，以及输出风格是否符合要求等定性指标。如何把提示集成进应用并传给模型，以及如何存储和维护提示，也在提示工程工具的功能范围内。

\begin{figure}[htbp]
\centering
\begin{tikzpicture}[
  font=\sffamily\small,
  operation/.style={draw=blue!50!black,rounded corners=3pt,
    align=center,minimum height=.77cm,text width=3.00cm,fill=white},
  decision/.style={diamond,draw=blue!50!black,aspect=2.85,inner sep=1pt,
    align=center,font=\sffamily\scriptsize,minimum width=2.6cm,
    minimum height=.72cm,fill=white},
  stage/.style={rounded corners=6pt,inner xsep=9pt,inner ysep=8pt},
  path/.style={-{Stealth[length=2mm]},draw=blue!50!black},
  answer/.style={fill=white,inner sep=1.5pt,font=\sffamily\scriptsize}]
\node[operation] (develop) at (0,0) {开发提示};
\node[operation] (debug) at (3.65,0) {调试、运行并评估};
\node[decision] (satisfied) at (3.65,-1.06) {结果满意？};
\node[operation] (modify) at (0,-1.06) {修改提示};
\node[operation] (dev) at (0,-2.68) {作为流水线步骤\\部署到开发环境};
\node[operation] (run) at (3.65,-2.68) {运行并测试流水线};
\node[decision] (qa) at (3.65,-3.74) {通过质量检查？};
\node[operation] (deploy) at (0,-5.39) {部署到生产环境};
\node[operation] (monitor) at (3.65,-5.39) {监控};
\begin{scope}[on background layer]
\node[stage,fill=blue!11,fit=(develop)(debug)(satisfied)(modify)] (design) {};
\node[stage,fill=yellow!21,fit=(dev)(run)(qa)] {};
\node[stage,fill=green!13,fit=(deploy)(monitor)] {};
\end{scope}
\node[above,anchor=south west,font=\sffamily\scriptsize] at (design.north west) {设计与开发};
\draw[path] (develop) -- (debug);
\draw[path] (debug) -- (satisfied);
\draw[path] (satisfied) -- node[answer,above]{否} (modify);
\draw[path] (satisfied.south) -- node[answer,above,sloped]{是} (dev.north east);
\draw[path] (dev) -- (run);
\draw[path] (run) -- (qa);
\draw[path] (qa.west) -- node[answer,above,sloped]{否} (dev.south east);
\draw[path] (qa.south) -- node[answer,above,sloped]{是} (deploy.north east);
\draw[path] (deploy) -- (monitor);
\draw[path] (monitor.south) to[out=-90,in=-90,looseness=.75] (-2.07,-5.37)
  to[out=90,in=180] (design.west);
\end{tikzpicture}
\caption{提示工程工作流概览。}
\label{fig:ch08-prompt-workflow}
\end{figure}


如第~3~章所述，提示设计既可以是简单的手工过程，也可以是复杂的自动优化过程。通常可从手工设计的提示模板入手。OpenAI Playground\footnote{\url{https://platform.openai.com/playground}}等工具对此很有帮助。以 OpenAI Playground 为例，用户可以方便地在 OpenAI 的不同大语言模型之间切换，考察模板能否推广到其他模型。用户还可考察提示设计与温度等模型超参数的相互作用。生成序列时，温度较低会使选择集中于概率最高的词元。提高温度则会增加词元选择的随机性。用户可以借此了解模型设置如何影响给定提示的响应质量（\refcite{ch08-ref-084}{Saravia, 2022}）。

另一种提示设计和优化工具是 Promptmetheus。\footnote{\url{https://promptmetheus.com/}}这是功能丰富的提示工程集成开发环境，支持编写、测试、优化及部署提示。原书当时的比较指出，OpenAI Playground 只支持与 OpenAI 的专有大语言模型交互，而 Promptmetheus 支持针对 70 多个专有和开源大语言模型开展开发。Promptmetheus 定位偏向开发工具，因此提供了更丰富的实验跟踪功能，也可与 LangChain 等开发工具集成。

\noindent\begin{minipage}{\linewidth}
\begin{lstlisting}[language=Python,caption={构造 \texttt{ChatPromptTemplate} 的示例。来源链接见正文。},label={lst:ch08-chat-template}]
from langchain_core.prompts import ChatPromptTemplate
chat_template = ChatPromptTemplate.from_messages(
    [
        ("system", "You are a helpful AI bot. Your name is {name}."),
        ("human", "Hello, how are you doing?"),
        ("ai", "I'm doing well, thanks!"),
        ("human", "{user_input}"),
    ]
)
messages = chat_template.format_messages(name="Bob", user_input="What is your name?")
\end{lstlisting}
\end{minipage}\par

LangChain 也是一种功能较复杂的工具，具备提示优化与管理功能。第~\ref{sec:ch08-frameworks}~节介绍过，它还提供多项功能，可用于构建包含多个组件的大语言模型应用。LangChain 的提示开发功能本身也很常用，为用户提供了多种便捷方法。例如，用户可借助 \texttt{ChatPromptTemplate} 以编程方式构造传给大语言模型的提示消息，见清单~\ref{lst:ch08-chat-template}。\footnote{示例来自 \url{https://python.langchain.com/docs/modules/model_io/prompts/quick_start/}。原书在此误引作清单 9.8。}这段代码生成一段多轮聊天对话，便于通过 LangChain 的多种方法之一传给模型。

设计提示时，还可以查阅已有提示库。提示库收录常用提示和设计或格式示例，通常稍作修改即可使用。例如，Prompt Hub\footnote{\url{https://www.promptingguide.ai/prompts}}覆盖多种任务和大语言模型交互模式，包含情感分类、SQL 查询生成、诗歌生成及实体抽取等任务的提示。库中未必有恰好符合需求的提示，但通常能找到有助于构造自己提示的示例，便于快速开始提示工程工作。

\subsection{评估与测试}
\label{sec:ch08-evaluation-testing}
部署任何机器学习应用之前，都有必要评估模型表现。在有真实答案的简单情形中，这项工作很直接。生成式应用往往没有唯一正确的输出，评估因而带有主观判断。一种常用方式是建立自动评估流程，由大语言模型本身或另一个能力更强的模型来判断输出质量。首先需要创建涵盖多种提示的测试用例，再让大语言模型为每个用例生成响应。输出可按多种标准评估，包括以下方面。
\begin{itemize}
\item 相关性
\item 全面性
\item 有据性
\item 幻觉
\item 所含个人身份信息（PII）
\item 情感倾向
\item 毒性
\end{itemize}

TruLens 等工具为多种常见评估需求提供内置提示，可以简化这一过程（\refcite{ch08-ref-079}{Reini et al., 2024}）。例如，只需一行函数调用即可取得某个响应的幻觉得分。如果某个应用需要专门的评估方法，也可以扩展这一框架。

% 原书第 8.7 节。
\section{推理}
\label{sec:ch08-inference}
为大语言模型应用设计服务系统时，需要就如何执行模型推理作出多项决定。大部分决定涉及成本、速度和模型表现三个方面的优化。700 亿参数的模型可以提供质量很高的输出，但其速度远低于 70 亿参数的模型，除非愿意在计算基础设施上投入大量资金。本节介绍在这些方面进行权衡的方法，以及因使用场景而异的若干重要因素。

\subsection{模型托管}
大语言模型应用的一项基本决策是模型本身托管在哪里。可以把选项分为以下三类。

\begin{itemize}
\item \textbf{向公开的第三方 API 发送推理请求。}这种方式快捷、直接，常用于构建演示和原型，以及迅速将新构想投入生产。它无需自行搭建或维护，开发人员无需深入了解大语言模型的原理就可学会使用这些 API。部署应用所需的工作和专业知识较少，可以节省不少费用。不过，API 本身的使用成本相对较高，推理请求量大时，很容易抵消所节省的开销。还有几个限制需要考虑。首先，这种方式几乎无法针对特定使用场景微调或以其他方式定制模型。其次，延迟没有强保证，还须考虑公开 API 常见的速率限制。最后，如第~6.4~节所述，这些 API 调用意味着应用中的数据会交给第三方。对许多组织来说，仅这一点就足以放弃该方案。尽管有多种缺点，OpenAI 当时最新的 GPT 模型只通过其 API 提供。如果采用 OpenAI 当时领先的模型能力可以为终端用户提供最大的价值，这些取舍可能是值得的。

\item \textbf{使用基础模型托管服务。}三家主要云计算提供商均有相应服务，例如 AWS Bedrock，可在私有、安全的环境中使用基础模型。出于多种原因，这种方式比公开 API 更容易扩展。首先，服务提供商为推理提供多项内置优化，账号所有者也可以在一定程度上控制分配给模型的 GPU 资源，从而在推理速度和计算开销之间找到合适的平衡。API 用户无法控制这两项因素。其次，如果大语言模型与应用的其他部分位于同一个云环境，许多网络问题可以减轻。推理请求和响应较少受带宽波动影响，也可以通过配置让模型与应用位于同一物理地点。延迟很重要时，把请求发送到可能远在地球另一端的 API 会带来问题。不过，模型托管服务也有限制。为了优化算力使用，服务所支持的基础模型和调优技术往往是预先规定的一组。优化还附带额外费用，对于内部已有 GPU 计算专长的组织，这笔开销未必划算。

\item \textbf{在自己的计算基础设施上托管模型。}如果 AWS Bedrock 等服务限制过多，可以自行搭建运行环境。这样可以灵活选择大语言模型，并针对应用需求精确优化，复杂度也高于前两种方案。NVIDIA 的 Triton 推理服务器\footnote{\url{https://github.com/triton-inference-server}}可以减少一部分工作。它支持灵活选择模型架构，并负责多项底层 GPU 优化。许多组织无需聘用或外包一支熟悉张量运算的专门团队，也能构建高度定制的推理系统。不过，推理需求达到很大规模时，这种投入可能带来成本优势。
\end{itemize}

\subsection{性能优化}
\label{sec:ch08-inference-performance}
在 CPU 上尝试过运行大语言模型的人，多半体会过 GPU 算力不足时的响应速度。算力成本较高，人们因此提出了多种方法，以便在不增加硬件的情况下更快地处理更多推理请求。

与推理速度有关的两个重要概念是延迟和吞吐量。延迟指处理一个请求并向应用发送响应所需的时间，吞吐量指大语言模型在给定时间内可以产生的输出量（\refcite{ch08-ref-001}{Agarwal et al., 2023}）。二者关系密切，但含义不同。例如，用户开始输入代码时，期待编程助手几乎立即给出建议。这类应用要求低延迟。再设想一个过滤垃圾邮件的服务。无论分类花半秒、几秒还是更久，用户很可能感觉不到区别，但吞吐量仍然重要。如果服务跟不上新邮件的涌入，积压就会越来越多，无法发挥预期作用。

\subsubsection{批处理}
\label{sec:ch08-batching}
优化批处理对于提高吞吐量很重要，见图~\ref{fig:ch08-batching}。把多个输入合入同一个矩阵，可以利用向量化带来的性能收益。一般来说，较大的批次比小批次效率更高。如果应用的大多数推理请求包含大量输入数据，批处理容易组织。若用户每次只发送一个输入，例如在对话应用中，则要把多个用户的输入合入同一批次，才能利用批处理。这称为动态批处理（dynamic batching）。它能提高吞吐量，但用户可能需要等待服务器积攒其他请求后才开始处理自己的请求，从而增大延迟。

连续批处理（continuous batching）是针对自回归大语言模型解决上述问题的方法。生成式文本模型迭代地产生词元，因而可以把新的用户输入加入正在处理的批次。若某个输出词元串已完成，即达到了长度上限或生成了停止词元，就会腾出批次中的一个输入位置。随后可以把队列中的下一个用户请求放入批次。这样，GPU 内存一旦腾出空间，系统就能开始处理新的请求，也能避免因批次小于合适规模而未充分利用 GPU。这种方法也能处理长度相差很大的输入，而不必为额外的填充词元付出代价。GPU 专长于大型矩阵运算，输入的大小和形状持续适合硬件架构时，其性能最高。

\sourcefigure{fig08_06.pdf}{标准批处理、动态批处理和连续批处理。传统推理架构主要由客户端组成批次。尤其在每个用户每次只发一个请求的应用中，服务器动态汇聚多个输入能更充分地利用 GPU，但凑齐批次后才能开始计算，因而会增加等待时间。连续批处理把新收到的输入加入正在执行的批次，以减少这种等待。}{fig:ch08-batching}

\subsubsection{键值缓存}
键值缓存（key-value caching）是另一种适用于自回归大语言模型的推理技术。每生成一个词元，就把它接在序列末尾，再将序列送入模型，以生成下一个词元。此前所有词元都与上一轮相同，因此其对应的键和值张量可缓存并反复使用，无需在每轮迭代中重新计算。这样只需为新词元生成键和值。\footnote{原文把可复用的量称为注意力权重。键值缓存存储的是此前词元的键和值，新词元仍需计算相对于这些键的注意力权重。}注意力机制是 Transformer 中时间复杂度最高的部分，通常也是架构中最大的性能瓶颈。减少这些计算可以显著提高推理速度。

\subsection{成本优化}
即便采用所有可用的推理加速技术，规模最大、能力最强的模型仍然需要价格较高的硬件。应用同时要求低延迟和高吞吐量时尤其如此。几乎总是值得考察较小的模型能否同样好地完成工作，或至少取得相近的表现。在某些使用场景中，限制客户选择的可能更多是预算，而非结果质量。降低成本与提高效用之间存在一系列可行取舍，生产应用都需要仔细分析，在其中找出合适的平衡点。

模型规模与结果质量之间还有一个可权衡的方面。采用较小模型节省的一部分费用可用来微调，以弥补能力差距。大模型的吸引力之一在于它包含丰富的知识，仅用提示工程和上下文学习，就能较好地完成广泛的任务。这很重要，因为即使采用 LoRA 等技术，微调大模型也很昂贵。选择较小的大语言模型后，微调更可行。Predibase 的业务建立在这一想法上。其主张是，多数应用适合使用较小的专用模型，并按各类推理请求的需要配备适配器。

% 原书第 8.8 节。
\section{大语言模型运维}
\label{sec:ch08-llmops}
随着人们对大语言模型的兴趣迅速增长，大语言模型运维（large language model operations，LLMOps）成为机器学习运维（MLOps）的一个自然分支。把大语言模型应用部署到生产环境时，许多问题与一般机器学习模型的部署问题相似。这里主要介绍与大语言模型有关的内容。有过其他模型运维经验的读者会熟悉其中不少方法。

当一种新思想迅速传播时，MLOps 和 LLMOps 往往成为流行词，具体涵盖什么也众说纷纭。这里采用较广义的理解，介绍能将机器学习能力部署到生产环境的多种工具与流程，包括数据和提示管理、迭代训练与工作流编排（\refcite{ch08-ref-064}{Oladele, 2024}）。多数方法旨在提高效率或降低风险，也有方法兼顾两者。这关系到能否从新的机器学习能力中获得价值。许多人曾把大量时间耗在模型上，给终端用户带来的好处很少，有时甚至有负面影响。

\subsection{LLMOps 工具与方法}
MLOps 及其后的 LLMOps 受到广泛重视，相应的解决方案市场迅速增长，出现了许多工具和产品。以下各节概述 LLMOps 生态，说明构建和部署大语言模型应用时会遇到哪些难题，以及常见的解决方式。图~\ref{fig:ch08-llmops}概述工具类别及其与生产应用的交互关系。

\sourcefigure{fig08_07.pdf}{一个基础聊天机器人及其 LLMOps 架构。提示模板经反复改进，并纳入版本控制。大语言模型也采用为该应用专门训练的适配器，因此需要跟踪实验并将训练好的组件发布到生产环境。提示模板和训练组件均通过 CI/CD 框架进入已部署的应用，生产指标则上报给监控服务。可以利用反馈进一步改进提示和训练数据。}{fig:ch08-llmops}

图的上部从提示工程、评估接到源码管理，从数据标注、训练流水线、实验跟踪接到模型注册库。下部依次展示用户界面、应用后端与推理服务器。后端用系统提示模板组织请求，可查询外部数据或服务，或直接响应用户。服务器负责请求路由、批处理，以及运行带适配器的预训练模型。示例中的用户查询订单，提供订单号 NX1234 后，收到商品缺货、预计于 2024 年 9 月 18 日发货的答复。

\subsubsection{实验跟踪}
把模型训练到生产部署所需的质量，需要反复迭代。特别是在开发新能力的早期，通常要进行大量试验。整理各轮结果并跟踪哪些模型表现最好，很快会变得繁重（Jakub \refcite{ch08-ref-032}{Czakon, 2024}）。日后可能还要找回其他细节，例如所用的具体数据集或所需的训练时间。许多实验跟踪工具可以满足这些需求。通常只需在训练流水线的代码中加入几条简单命令，相关信息便会自动记录。此类工具一般有功能完善的用户界面，提供多种指标和可视化方式，以便保存和分析实验。多人参与项目、希望查看彼此的工作时，这一点对团队协作特别有用。

模型注册库常与实验跟踪工具配合使用。模型注册库是保存通过实验跟踪工具创建的各个模型的仓库，也可直接上传未经实验跟踪的模型。模型注册库通常允许为模型添加自定义标签。下游流程可以据此自动识别哪些模型已可部署，或触发其他工作流。

\subsubsection{版本控制}
源码管理是软件项目中广泛采用的做法，同样适用于大语言模型。使用 Git 等代码仓库管理大语言模型的训练与评估代码，方式大体与管理其他代码库相同。以下还有一些仅靠代码仓库难以解决的版本管理需求。

\begin{itemize}
\item \textbf{模型版本管理。}构建大语言模型及其应用时，要多次迭代训练与调优。必须知道投入生产的是哪一版模型，并能找到生成该版本时用的确切代码和数据。否则，模型表现异常时很难调试并找出原因。采用实验跟踪工具和模型注册库后，这些工作大多已有工具支持。即便开发流程还没有那么成熟，无法使用先进工具，也必须以某种形式管理模型版本。
\item \textbf{数据版本管理。}训练数据往往是项目中变化最大的部分。常常要花大量时间确定最合适的数据类型，而用于编写代码的时间少得多。如果没有数据版本记录，模型就无法回滚至先前状态，相当于丢掉了以往的工作历史。
\item \textbf{提示版本管理。}提示和提示模板也是大语言模型系统的重要部分，可能在应用生命周期中发生很大变化。提示经常与应用代码一起存储，但这种做法未必始终合适。提示模板的行为通常类似概率模型，其验证方法因而与测试确定性代码的方式大不相同。单独管理提示可以简化开发，使团队能快速迭代提示，而无需每次都发布并部署较大组件的新版本。
\end{itemize}

\subsubsection{部署}
成熟软件开发生命周期中，许多常用于持续集成和持续交付或部署（CI/CD）的工具也适合部署大语言模型功能。这些流程旨在自动构建和测试待发布的新组件，并落实以下做法。

\begin{itemize}
\item 每次发布新版本都必须运行并通过同一组测试，以减少功能退化的可能。
\item 在通常模拟生产环境的集中环境中验证所有组件，统一各开发人员所用的验证条件。
\item 构建和发布过程应完全自动化且可重复，避免人工误操作使部署版本与通过测试的版本不同。
\item 组件投入生产后若表现未达要求，应能够迅速回滚至先前版本。
\end{itemize}

这种系统可以自动评估模型和报告指标，从而阻止表现欠佳的模型被误部署到生产环境。

\subsubsection{监控}
有许多工具可以监控已部署模型的表现，通常采用实时评估并汇总相关指标的方式。如果模型一开始就不符合预期，或其表现随时间发生变化，系统可以发出告警。例如，输出总体上比过去更短或更长，可能说明用户行为已经改变，也可能有其他原因。此时值得调查是否要相应调整模型或提示模板等组件。监控工具除了观察模型表现，还可防范模型幻觉和提示注入攻击等问题（\refcite{ch08-ref-004}{Ama, 2023}）。

可观测性与监控密切相关，两个术语有重叠，常被互换使用，市面上多数 LLMOps 方案也同时提供这两类功能。粗略区分的话，监控根据系统的汇总表现找出问题或有待改善的方面，可观测性则提供机制来进一步调查。例如，监控工具可能提示某天大语言模型的响应比平时慢。缺乏适当的可观测性时，可能难以找出原因。有了相应工具，就可以找出个别输入，逐步跟踪它在系统中的处理过程，定位瓶颈或故障。

\subsection{缩短迭代周期}
前面的章节讨论过人类反馈对于提高模型质量的重要性。人类反馈对基础模型的对齐和指令调优有价值，用户反馈对应用同样有价值。团队越快收集生产系统的结果，并用这些结果构建和发布更新，就越能为客户提供价值。改进这项工作流程，还可以迅速应对无法预见的问题。

另一个应缩短模型部署周期的理由是，用户反馈常常无可替代（\refcite{ch08-ref-012}{Burgess, 2021}）。数据科学家和开发人员做大量测试，也许能修正许多模型缺陷。不过，在处理真实用户输入之前，无法保证会得到预期结果。很多团队曾在新技术上投入数月乃至数年的努力，最终却在推出产品时受挫。避免这种局面的一种方式，是尽快推出最早可用的模型，同时准备好基础设施，以便取得新数据后可以迅速更新或回滚。ChatGPT 发布不久就出现了这种情况。谷歌和微软等主要竞争对手感到有必要尽快推出自己的聊天功能。好奇的用户很快开始攻击这些模型，从奇怪的响应中找乐趣，也暴露了模型的一些问题（\refcite{ch08-ref-065}{Orf, 2023}）。这些公司迅速处理早期问题，产品口碑最终没有受到太大影响。

\subsubsection{自动重新训练}
在某些情况下，可以直接用生产环境输入中的样例更新模型。例如，许多平台允许用户标记喜欢或不喜欢的内容。这类反馈可直接并入标注数据集，供下一轮训练使用。如果生产模型已较成熟，获取新数据、运行训练流水线、验证结果和部署新版本都可以全自动进行。需要快速适应新趋势的应用值得争取做到这一点，但实现起来并不容易。控制措施不足时，训练不充分的模型可能进入生产环境。应尽早考虑相应风险及降低这些风险所需的工作，再决定是否采用完全自动化的训练。达到这一成熟程度往往需要长期投入建设 LLMOps 能力。

\subsubsection{人在回路中的流程}
自动重新训练不可行或没有必要时，还有其他提高效率的方法。数据需要由人类专家审查或标注时，可采用多种标注工具。一些工具还可按需运用主动学习或半监督技术加快标注。前面讨论过，从人类反馈中进行强化学习（RLHF）显著提升了大语言模型的表现，所需的昂贵人力投入因此有了回报。收集或生成针对模型弱点的高质量训练样例，影响很大。样例的质量通常随人工投入增加。组织若无法为标注工作配备足够人员，也可将其外包，许多公司专门提供这类服务。

较新的数据或质量更可靠的数据往往会带来改进，可以迅速得到适于生产的模型，也有例外。模型架构可能是限制因素，这就需要进一步探索和研究，解决应用的弱点。实验跟踪和版本管理在这里很有帮助，特别是多人从不同角度研究同一个问题时。人工工作越多，越容易发生结果放错位置、数据集被改动等错误。还可以建立一个有代表性的黄金数据集供评估使用，并在整个迭代周期中保持固定，以便比较不同实验。

\subsection{风险管理}
用大语言模型或任何机器学习模型辅助完成传统上由人承担的任务，本身就涉及许多风险。设想一家公司被发现利用人工智能，以歧视性理由拒绝向某些人提供住房或就业机会，这可能造成严重影响。在这类环境中部署大语言模型时，需要特别认真地评估并尽量减少风险。

\begin{keypoint}
\textbf{风险评估与缓解}

根据大语言模型应用的功能进行充分的风险评估，是投入生产的重要一步。应考虑以下风险的影响。
\begin{itemize}
\item \textbf{侵犯隐私。}在正常使用过程中，应用是否会处理可能被意外展示给用户的受保护数据？
\item \textbf{安全漏洞。}应用设计和功能是否有可能导致数据丢失的漏洞？除了传统应用安全，还要考虑大语言模型所特有的提示注入攻击等风险。
\item \textbf{偏见和有害性。}应了解模型生成带有偏见或有害输出的风险。有些应用领域有明确的法律或监管规定，在这些领域尤其需要谨慎。
\end{itemize}
完成全面的风险评估后，可能需要采取缓解措施，包括以下几项。
\begin{itemize}
\item \textbf{人在回路中的监督。}许多请求量小、敏感度高或风险较大的使用场景，可能需要让人类专家参与。应用用户应能够在采取进一步行动之前核查大语言模型的输出。
\item \textbf{内容过滤。}应用可能需要实施内容过滤策略。可通过对齐调优或生成后过滤实现，例如采用关键词、模式匹配或分类方法，只将适宜的输出提供给用户。
\item \textbf{访问控制。}传统应用开发常有访问控制。在应用中使用大语言模型，可能还需要考虑其他角色权限。例如，对提示模板的操作权限可能只应授予担任提示工程师的开发人员。
\item \textbf{持续监控。}生产应用应记录正常的用户行为与交互情况，并持续监控。对大语言模型应用而言，了解用量、输出和表现有助于尽早发现问题。
\end{itemize}
\end{keypoint}

\subsubsection{模型治理}
前面已经看到，大语言模型仍面临若干问题，偏见只是其中之一。组织出于合理考虑，在受法律或监管约束的场景中使用机器学习时通常格外谨慎，尤其是医疗、金融和法律的多数领域。开发这类应用的技术人员应事先考虑要满足哪些要求，才能让利益相关者相信大语言模型的收益超过风险。许多组织制定了标准，要求对生产模型充分验证并记录。模型可解释性也可能至关重要。达到这些标准的工作量容易被低估。如果项目日程没有为此留出时间，就可能发生延误并增加成本。

模型卡（model card）是较常用的模型治理手段（\refcite{ch08-ref-059}{Mitchell et al., 2019}）。这项标准鼓励模型开发人员公开相关信息，降低模型被用于预期之外用途的风险。信息应以便于任何模型使用者查阅的方式呈现。不同模型适用的条目不尽相同。理想情况下，模型卡应说明推荐用途、已知限制和训练数据中的潜在偏见，也可以提供训练过程和不同基准上的评估指标。模型卡通常简明，技术细节较少。

\subsubsection{数据治理}
大语言模型快速兴起的一个重要原因是从网上抓取了大量数据。数据集迅速扩大后，整理或过滤特定类别的数据越来越难。大语言模型进入主流后，越来越多的人意识到自己的数据被用于未同意、甚至从未想过的用途，由此产生了矛盾。开发人员必须认真对待这些顾虑，以保护所在组织免受法律纠纷。

从社交媒体或论坛等网站提取数据时，可以先查阅相关条款。有些数据的使用规定可能不明确，也可能眼下似乎可以使用，日后却受到审查。因此，宜记录数据来源，包括每个数据集来自哪里、哪些模型使用过这些来源的数据。如果某些数据的使用因隐私、版权或其他问题而受到质疑，就可以减少损失，从存储中清除有关数据集，之后用排除这些数据的训练集训练模型。

% 原书第 8.9 节。
\section{教程：准备将实验模型部署到生产环境}
\label{sec:ch08-tutorial}

\subsection{概述}
本教程回顾第~4~章教程得到的实验模型。这次考察准备把其中一个模型部署到生产应用时可以采取的措施，演示本章讨论的若干工具和技术。整个操作仍在 Colab 笔记本环境中完成，因为许多读者可能希望避免实际部署生产级推理能力的开销。

\begin{keypoint}
\textbf{学习目标}
\begin{itemize}
\item 试用一种开源评估工具和一种开源监控工具。
\item 探索这些工具提供的功能及其用途。
\item 观察该过程是否揭示了可能影响模型能否部署到生产环境的新特征。
\end{itemize}
\end{keypoint}

本节为教程的简版，完整版见
\url{https://github.com/springer-llms-deep-dive/llms-deep-dive-tutorials}。

\subsection{实验设计}
本练习聚焦于把大语言模型能力从构想推进到生产时值得考虑的若干因素。设想我们在第~4~章的教程中得到两个候选模型。通过比较两者的长处和短处，确定哪一个更适合设想中的应用，并考察能否消除计算瓶颈来控制推理成本。然后考虑所选模型部署后的长期影响，展示如何检查它仍能完成原定任务，且没有意外后果。

首先进行模型评估，以便在投入运行前充分考察模型行为。第~4~章手工向 GPT-4 提供评分规则来评估模型。这里采取相似的方法，改用开源工具 TruLens（\refcite{ch08-ref-079}{Reini et al., 2024}）。它提供可扩展的评估框架及用于比较模型指标的仪表板。市面上有若干类似方案，TruLens 的一项优点是免费。

接着简要考察模型的推理速度。实际应用中，可以在不同 GPU 架构上进行基准测试，并考虑每种架构可采用的优化，以了解运行某个模型的成本。在本练习中，仅观察模型在 Colab GPU 上的运行情况。

教程的最后假设模型已经在生产中运行了一段时间，希望查看它的行为是否仍符合预期，以及系统中是否发生了可能影响模型表现的变化。为演示这一点，有意调整一部分测试数据，使用户提示呈现逐渐变长的趋势。这一环节采用 WhyLabs 提供的另一款免费开源工具 LangKit。

\subsection{结果与分析}
首先用 TWEETSUMM 测试集中的一小部分样本演示 \texttt{trulens\_eval} 库。TruLens 使用反馈函数进行评估，既可用内置函数，也可用自定义函数。本练习选择内置的连贯性与简洁性反馈函数。TruLens 在内部封装了 OpenAI、LangChain 等 API，开发人员可选择希望采用的服务提供方。简洁性等指标通过提示模板获得。

\noindent\begin{minipage}{\linewidth}
\begin{lstlisting}[language=Python,caption={用于 TruLens 评估的系统提示模板示例。},label={lst:ch08-trulens-prompt}]
f"""{supported_criteria['conciseness ']} Respond only as a
number from 0 to 10 where 0 is the least concise and 10 is
the most concise."""
\end{lstlisting}
\end{minipage}\par

将 DistilGPT2 和 Llama-2 LoRA 模型用于测试样本，得到表~\ref{tab:ch08-trulens}所示的平均分。TruLens 的各项指标均按 0 到 1 评分。如预期，较大的 Llama-2 模型各项得分都较高。连贯性与简洁性得分看起来较合理，摘要质量得分可能偏低，DistilGPT2 尤其如此。此前的教程中，这些模型的表现看起来不错。这很可能部分是因为本练习自定义评估时，没有花太多时间设计提示模板。TruLens 的连贯性和简洁性评估则建立在经开发人员用一组测试用例检验过的提示模板上。这一示例说明评估并不容易，经过检验的现有方案因而很有价值。

\begin{table}[htbp]
\centering\small
\caption{用 TruLens 评估两个候选模型的结果。连贯性和简洁性是工具内置指标，摘要质量是我们提供的自定义评估指标。}
\label{tab:ch08-trulens}
\begin{tabular}{@{}lrrr@{}}
\toprule
模型 & 平均连贯性 & 平均简洁性 & 平均摘要质量\\
\midrule
DistilGPT2-finetuned & 0.66 & 0.80 & 0.29\\
Llama-2 LoRA & 0.80 & 0.83 & 0.60\\
\bottomrule
\end{tabular}
\end{table}

采用统一的评估格式，并尽量复用现有提示，有若干好处。首先，可以节省设计评估方法的时间。定义这类抽象接口也可以使 LLMOps 系统的不同环节更容易实现自动化。例如，TruLens 可以接入应用，实时评估经过系统的用户输入和模型输出，以便立即提供反馈。本例未模拟这种操作。

然后使用另一款免费的 LLMOps 工具 LangKit。它是 WhyLabs 监控与可观测性软件套件的一部分。这里考察它分析提示与响应随时间变化趋势的能力。为模拟变化，我们建立并比较两个独立的数据批次，即数据概况（profile）。把数据分为两个较小的集合，分别包含较长与较短的输入，使数据概况之间出现差异。随后连接 WhyLabs 仪表板，详细查看各项指标，见图~\ref{fig:ch08-whylabs}。

\sourcefigure{fig08_08.pdf}{WhyLabs 监控仪表板。查看选定指标，了解模拟的提示数据漂移会产生怎样的影响。}{fig:ch08-whylabs}

将 TruLens 和 LangKit 都用于 TWEETSUMM 模型和数据后，可以看到两款工具的功能有所重叠。不过，它们的实现方式很不相同，也各有长处。TruLens 偏重评估，LangKit 则偏重日志记录与监控。视应用而定，可以两款都用，也可以从中择一。它们只是众多 LLMOps 方案中的两种，通常还需要调研才能找到合适的做法。

\subsection{结论}
将大语言模型应用投入生产涉及大量工作，远超本简短练习能覆盖的范围。练习演示了可以降低这一过程复杂度的若干工具。判断模型能否投入生产要考虑许多方面，TruLens 和 LangKit 等工具的开发者已经为其中许多问题提供了可重复使用的方案。基于这些工具构建工作流，可以缩短应用达到较成熟状态所需的时间。

% PDF 392--400。
\begin{chapterreferences}
\refentry{ch08-ref-001}{Megha Agarwal, Asfandyar Qureshi, Nikhil Sardana, Linden Li, Julian Quevedo, and Daya Khudia. Llm inference performance engineering: Best practices, 10 2023. URL \url{https://www.databricks.com/blog/llm-inference-performance-engineering-best-practices}.}
\refentry{ch08-ref-002}{Amey Agrawal, Nitin Kedia, Ashish Panwar, Jayashree Mohan, Nipun Kwatra, Bhargav S. Gulavani, Alexey Tumanov, and Ramachandran Ramjee. Taming throughput-latency tradeoff in llm inference with sarathi-serve, 2024.}
\refentry{ch08-ref-003}{Norah Alzahrani et al. When benchmarks are targets: Revealing the sensitivity of large language model leaderboards, 2024.}
\refentry{ch08-ref-004}{Emeka Boris Ama. Llm monitoring: The beginner’s guide, 11 2023. URL \url{https://www.lakera.ai/blog/llm-monitoring}.}
\refentry{ch08-ref-005}{Stéphane Aroca-Ouellette, Cory Paik, Alessandro Roncone, and Katharina Kann. Prost: Physical reasoning of objects through space and time. \textit{arXiv preprint} \textit{arXiv:2106.03634}, 2021.}
\refentry{ch08-ref-006}{Jacob Austin, Augustus Odena, Maxwell Nye, Maarten Bosma, Henryk Michalewski, David Dohan, Ellen Jiang, Carrie Cai, Michael Terry, Quoc Le, et al. Program synthesis with large language models. \textit{arXiv preprint} \textit{arXiv:2108.07732}, 2021.}
\refentry{ch08-ref-007}{Gad Benram. Top tools for prompt engineering?, 2023. URL \url{https://www.tensorops.ai/post/top-tools-for-prompt-engineering}.}
\refentry{ch08-ref-008}{Yonatan Bisk, Rowan Zellers, Jianfeng Gao, Yejin Choi, et al. Piqa: Reasoning about physical commonsense in natural language. In \textit{Proceedings of the AAAI conference on artificial intelligence}, volume 34, pages 7432–7439, 2020.}
\refentry{ch08-ref-009}{Su Lin Blodgett, Lisa Green, and Brendan O’Connor. Demographic dialectal variation in social media: A case study of african-american english. \textit{arXiv preprint} \textit{arXiv:1608.08868}, 2016.}
\refentry{ch08-ref-010}{Ondrej Bojar, Rajen Chatterjee, Christian Federmann, Yvette Graham, Barry Haddow, Matthias Huck, Antonio Jimeno Yepes, Philipp Koehn, Varvara Logacheva, Christof Monz, et al. Findings of the 2016 conference on machine translation (wmt16). In \textit{First conference on machine translation}, pages 131–198. Association for Computational Linguistics, 2016.}
\refentry{ch08-ref-011}{Jordan Boyd-Graber, Brianna Satinoff, He He, and Hal Daumé III. Besting the quiz master: Crowdsourcing incremental classification games. In \textit{Proceedings of the} \textit{2012 joint conference on empirical methods in natural language processing and} \textit{computational natural language learning}, pages 1290–1301, 2012.}
\refentry{ch08-ref-012}{Jordan Burgess. What is human-in-the-loop ai?, 11 2021. URL \url{https://humanloop.com/blog/human-in-the-loop-ai}.}
\refentry{ch08-ref-013}{Yingshan Chang, Mridu Narang, Hisami Suzuki, Guihong Cao, Jianfeng Gao, and Yonatan Bisk. Webqa: Multihop and multimodal qa. In \textit{Proceedings of} \textit{the IEEE/CVF Conference on Computer Vision and Pattern Recognition}, pages 16495–16504, 2022.}
\refentry{ch08-ref-014}{Mark Chen et al. Evaluating large language models trained on code, 2021.}
\refentry{ch08-ref-015}{Eunsol Choi, He He, Mohit Iyyer, Mark Yatskar, Wen-tau Yih, Yejin Choi, Percy Liang, and Luke Zettlemoyer. Quac: Question answering in context. \textit{arXiv} \textit{preprint arXiv:1808.07036}, 2018.}
\refentry{ch08-ref-016}{Christopher Clark, Kenton Lee, Ming-Wei Chang, Tom Kwiatkowski, Michael Collins, and Kristina Toutanova. BoolQ: Exploring the surprising difficulty of natural yes/no questions. In \textit{Proceedings of NAACL-HLT 2019}, 2019.}
\refentry{ch08-ref-017}{Jonathan H Clark, Eunsol Choi, Michael Collins, Dan Garrette, Tom Kwiatkowski, Vitaly Nikolaev, and Jennimaria Palomaki. Tydi qa: A benchmark for information-seeking question answering in typologically diverse languages. \textit{Transactions of} \textit{the Association for Computational Linguistics}, 8:454–470, 2020.}
\refentry{ch08-ref-018}{Peter Clark, Isaac Cowhey, Oren Etzioni, Tushar Khot, Ashish Sabharwal, Carissa Schoenick, and Oyvind Tafjord. Think you have solved question answering? try arc, the ai2 reasoning challenge. \textit{arXiv preprint arXiv:1803.05457}, 2018.}
\refentry{ch08-ref-019}{Karl Cobbe, Vineet Kosaraju, Mohammad Bavarian, Mark Chen, Heewoo Jun, Lukasz Kaiser, Matthias Plappert, Jerry Tworek, Jacob Hilton, Reiichiro Nakano, et al. Training verifiers to solve math word problems. \textit{arXiv preprint} \textit{arXiv:2110.14168}, 2021.}
\refentry{ch08-ref-020}{Alexis Conneau, Guillaume Lample, Ruty Rinott, Adina Williams, Samuel R Bowman, Holger Schwenk, and Veselin Stoyanov. Xnli: Evaluating cross-lingual sentence representations. \textit{arXiv preprint arXiv:1809.05053}, 2018.}
\refentry{ch08-ref-021}{Ido Dagan, Oren Glickman, and Bernardo Magnini. The pascal recognising textual entailment challenge. In \textit{Machine learning challenges workshop}, pages 177–190. Springer, 2005.}
\refentry{ch08-ref-022}{Marie-Catherine De Marneffe, Mandy Simons, and Judith Tonhauser. The commitmentbank: Investigating projection in naturally occurring discourse. In \textit{proceedings of Sinn und Bedeutung}, volume 23, pages 107–124, 2019.}
\refentry{ch08-ref-023}{Emily Dinan, Stephen Roller, Kurt Shuster, Angela Fan, Michael Auli, and Jason Weston. Wizard of wikipedia: Knowledge-powered conversational agents. \textit{arXiv} \textit{preprint arXiv:1811.01241}, 2018.}
\refentry{ch08-ref-024}{Emily Dinan, Varvara Logacheva, Valentin Malykh, Alexander Miller, Kurt Shuster, Jack Urbanek, Douwe Kiela, Arthur Szlam, Iulian Serban, Ryan Lowe, et al. The second conversational intelligence challenge (convai2). In \textit{The NeurIPS’18 Competition: From Machine Learning to Intelligent Conversations}, pages 187–208. Springer, 2020.}
\refentry{ch08-ref-025}{Dheeru Dua, Yizhong Wang, Pradeep Dasigi, Gabriel Stanovsky, Sameer Singh, and Matt Gardner. Drop: A reading comprehension benchmark requiring discrete reasoning over paragraphs. \textit{arXiv preprint arXiv:1903.00161}, 2019.}
\refentry{ch08-ref-026}{Thiago Castro Ferreira, Claire Gardent, Nikolai Ilinykh, Chris Van Der Lee, Simon Mille, Diego Moussallem, and Anastasia Shimorina. The 2020 bilingual, bidirectional webnlg+ shared task overview and evaluation results (webnlg+ 2020). In \textit{Proceedings of the 3rd International Workshop on Natural Language Generation from the Semantic Web (WebNLG+)}, 2020.}
\refentry{ch08-ref-027}{Sarah Gao and Andrew Kean Gao. On the origin of llms: An evolutionary tree and graph for 15,821 large language models, 2023.}
\refentry{ch08-ref-028}{Mor Geva, Daniel Khashabi, Elad Segal, Tushar Khot, Dan Roth, and Jonathan Berant. Did aristotle use a laptop? a question answering benchmark with implicit reasoning strategies. \textit{Transactions of the Association for Computational Linguistics}, 9:346–361, 2021.}
\refentry{ch08-ref-029}{Dan Hendrycks, Collin Burns, Steven Basart, Andy Zou, Mantas Mazeika, Dawn Song, and Jacob Steinhardt. Measuring massive multitask language understanding. \textit{arXiv preprint arXiv:2009.03300}, 2020.}
\refentry{ch08-ref-030}{Dan Hendrycks, Steven Basart, Saurav Kadavath, Mantas Mazeika, Akul Arora, Ethan Guo, Collin Burns, Samir Puranik, Horace He, Dawn Song, et al. Measuring coding challenge competence with apps. \textit{arXiv preprint arXiv:2105.09938}, 2021.}
\refentry{ch08-ref-031}{Jordan Hoffmann et al. Training compute-optimal large language models, 2022.}
\refentry{ch08-ref-032}{Kilian Kluge Jakub Czakon. Ml experiment tracking: What it is, why it matters, and how to implement it, 4 2024. URL \url{https://neptune.ai/blog/ml-experiment-tracking}.}
\refentry{ch08-ref-033}{Jeff Johnson, Matthijs Douze, and Hervé Jégou. Billion-scale similarity search with gpus, 2017.}
\refentry{ch08-ref-034}{Mandar Joshi, Eunsol Choi, Daniel S. Weld, and Luke Zettlemoyer. Triviaqa: A large scale distantly supervised challenge dataset for reading comprehension, 2017.}
\refentry{ch08-ref-035}{Daniel Khashabi, Snigdha Chaturvedi, Michael Roth, Shyam Upadhyay, and Dan Roth. Looking beyond the surface: A challenge set for reading comprehension over multiple sentences. In \textit{Proceedings of the 2018 Conference of the North American Chapter of the Association for Computational Linguistics: Human Language} \textit{Technologies, Volume 1 (Long Papers)}, pages 252–262, 2018.}
\refentry{ch08-ref-036}{Rik Koncel-Kedziorski, Subhro Roy, Aida Amini, Nate Kushman, and Hannaneh Hajishirzi. Mawps: A math word problem repository. In \textit{Proceedings of the 2016} \textit{conference of the north american chapter of the association for computational} \textit{linguistics: human language technologies}, pages 1152–1157, 2016.}
\refentry{ch08-ref-037}{Sumith Kulal, Panupong Pasupat, Kartik Chandra, Mina Lee, Oded Padon, Alex Aiken, and Percy S Liang. Spoc: Search-based pseudocode to code. \textit{Advances in} \textit{Neural Information Processing Systems}, 32, 2019.}
\refentry{ch08-ref-038}{Matt Kusner, Yu Sun, Nicholas Kolkin, and Kilian Weinberger. From word embeddings to document distances. In Francis Bach and David Blei, editors, \textit{Proceedings of the 32nd International Conference on Machine Learning}, volume 37 of \textit{Proceedings of Machine Learning Research}, pages 957–966, Lille, France, 07–09 Jul 2015. PMLR. URL \url{https://proceedings.mlr.press/v37/kusnerb15.html}.}
\refentry{ch08-ref-039}{Tom Kwiatkowski, Jennimaria Palomaki, Olivia Redfield, Michael Collins, Ankur Parikh, Chris Alberti, Danielle Epstein, Illia Polosukhin, Jacob Devlin, Kenton Lee, et al. Natural questions: a benchmark for question answering research. \textit{Transactions of the Association for Computational Linguistics}, 7:453–466, 2019.}
\refentry{ch08-ref-040}{Yanis Labrak, Adrien Bazoge, Emmanuel Morin, Pierre-Antoine Gourraud, Mickael Rouvier, and Richard Dufour. Biomistral: A collection of open-source pretrained large language models for medical domains, 2024.}
\refentry{ch08-ref-041}{Guokun Lai, Qizhe Xie, Hanxiao Liu, Yiming Yang, and Eduard Hovy. Race: Large-scale reading comprehension dataset from examinations, 2017.}
\refentry{ch08-ref-042}{Yuhang Lai, Chengxi Li, Yiming Wang, Tianyi Zhang, Ruiqi Zhong, Luke Zettlemoyer, Wen-tau Yih, Daniel Fried, Sida Wang, and Tao Yu. Ds-1000: A natural and reliable benchmark for data science code generation. In \textit{International Conference on Machine Learning}, pages 18319–18345. PMLR, 2023.}
\refentry{ch08-ref-043}{Huong T Le, Dung T Cao, Trung H Bui, Long T Luong, and Huy Q Nguyen. Improve quora question pair dataset for question similarity task. In \textit{2021 RIVF International Conference on Computing and Communication Technologies (RIVF)}, pages 1–5. IEEE, 2021.}
\refentry{ch08-ref-044}{Hector Levesque, Ernest Davis, and Leora Morgenstern. The winograd schema challenge. In \textit{Thirteenth international conference on the principles of knowledge representation and reasoning}, 2012.}
\refentry{ch08-ref-045}{Patrick Lewis, Barlas Oğuz, Ruty Rinott, Sebastian Riedel, and Holger Schwenk. Mlqa: Evaluating cross-lingual extractive question answering. \textit{arXiv preprint} \textit{arXiv:1910.07475}, 2019.}
\refentry{ch08-ref-046}{Peng Li, Wei Li, Zhengyan He, Xuguang Wang, Ying Cao, Jie Zhou, and Wei Xu. Dataset and neural recurrent sequence labeling model for open-domain factoid question answering. \textit{arXiv preprint arXiv:1607.06275}, 2016.}
\refentry{ch08-ref-047}{Yinheng Li, Shaofei Wang, Han Ding, and Hang Chen. Large language models in finance: A survey, 2023.}
\refentry{ch08-ref-048}{Yujia Li, David Choi, Junyoung Chung, Nate Kushman, Julian Schrittwieser, Rémi Leblond, Tom Eccles, James Keeling, Felix Gimeno, Agustin Dal Lago, et al. Competition-level code generation with alphacode. \textit{Science}, 378(6624):1092– 1097, 2022.}
\refentry{ch08-ref-049}{Percy Liang, Rishi Bommasani, Tony Lee, Dimitris Tsipras, Dilara Soylu, Michihiro Yasunaga, Yian Zhang, Deepak Narayanan, Yuhuai Wu, Ananya Kumar, et al. Holistic evaluation of language models. \textit{arXiv preprint arXiv:2211.09110}, 2022.}
\refentry{ch08-ref-050}{Wang Ling, Dani Yogatama, Chris Dyer, and Phil Blunsom. Program induction by rationale generation: Learning to solve and explain algebraic word problems. \textit{arXiv preprint arXiv:1705.04146}, 2017.}
\refentry{ch08-ref-051}{Bang Liu, Di Niu, Haojie Wei, Jinghong Lin, Yancheng He, Kunfeng Lai, and Yu Xu. Matching article pairs with graphical decomposition and convolutions. \textit{arXiv} \textit{preprint arXiv:1802.07459}, 2018.}
\refentry{ch08-ref-052}{Jian Liu, Leyang Cui, Hanmeng Liu, Dandan Huang, Yile Wang, and Yue Zhang. Logiqa: A challenge dataset for machine reading comprehension with logical reasoning. \textit{arXiv preprint arXiv:2007.08124}, 2020.}
\refentry{ch08-ref-053}{Chung Kwan Lo. What is the impact of chatgpt on education? a rapid review of the literature. \textit{Education Sciences}, 13(4):410, 2023.}
\refentry{ch08-ref-054}{Barrault Loïc, Biesialska Magdalena, Bojar Ondřej, Federmann Christian, Graham Yvette, Grundkiewicz Roman, Haddow Barry, Huck Matthias, Joanis Eric, Kocmi Tom, et al. Findings of the 2020 conference on machine translation (wmt20). In \textit{Proceedings of the Fifth Conference on Machine Translation}, pages 1–55. Association for Computational Linguistics, 2020.}
\refentry{ch08-ref-055}{R Thomas McCoy, Ellie Pavlick, and Tal Linzen. Right for the wrong reasons: Diagnosing syntactic heuristics in natural language inference. \textit{arXiv preprint} \textit{arXiv:1902.01007}, 2019.}
\refentry{ch08-ref-056}{Stephen Merity, Caiming Xiong, James Bradbury, and Richard Socher. Pointer sentinel mixture models. \textit{arXiv preprint arXiv:1609.07843}, 2016.}
\refentry{ch08-ref-057}{Shen-Yun Miao, Chao-Chun Liang, and Keh-Yih Su. A diverse corpus for evaluating and developing english math word problem solvers. \textit{arXiv preprint} \textit{arXiv:2106.15772}, 2021.}
\refentry{ch08-ref-058}{Todor Mihaylov, Peter Clark, Tushar Khot, and Ashish Sabharwal. Can a suit of armor conduct electricity? a new dataset for open book question answering. \textit{arXiv} \textit{preprint arXiv:1809.02789}, 2018.}
\refentry{ch08-ref-059}{Margaret Mitchell, Simone Wu, Andrew Zaldivar, Parker Barnes, Lucy Vasserman, Ben Hutchinson, Elena Spitzer, Inioluwa Deborah Raji, and Timnit Gebru. Model cards for model reporting. In \textit{Proceedings of the Conference on Fairness, Accountability, and Transparency}, FAT* ’19. ACM, January 2019. doi: 10.1145/3287560. 3287596. URL \url{http://dx.doi.org/10.1145/3287560.3287596}.}
\refentry{ch08-ref-060}{Nasrin Mostafazadeh, Nathanael Chambers, Xiaodong He, Devi Parikh, Dhruv Batra, Lucy Vanderwende, Pushmeet Kohli, and James Allen. A corpus and evaluation framework for deeper understanding of commonsense stories. \textit{arXiv preprint} \textit{arXiv:1604.01696}, 2016.}
\refentry{ch08-ref-061}{Shashi Narayan, Shay B Cohen, and Mirella Lapata. Don’t give me the details, just the summary! topic-aware convolutional neural networks for extreme summarization. \textit{arXiv preprint arXiv:1808.08745}, 2018.}
\refentry{ch08-ref-062}{Yixin Nie, Adina Williams, Emily Dinan, Mohit Bansal, Jason Weston, and Douwe Kiela. Adversarial nli: A new benchmark for natural language understanding. \textit{arXiv preprint arXiv:1910.14599}, 2019.}
\refentry{ch08-ref-063}{Jekaterina Novikova, Ondřej Dušek, and Verena Rieser. The e2e dataset: New challenges for end-to-end generation. \textit{arXiv preprint arXiv:1706.09254}, 2017.}
\refentry{ch08-ref-064}{Stephen Oladele. Llmops: What it is, why it matters, and how to implement it, 3 2024. URL \url{https://neptune.ai/blog/llmops}.}
\refentry{ch08-ref-065}{Darren Orf. Microsoft has lobotomized the ai that went rogue, 2 2023. URL \url{https://www.popularmechanics.com/technology/robots/a43017405/microsoft-bing-ai-chatbot-problems/}.}
\refentry{ch08-ref-066}{Denis Paperno, Germán Kruszewski, Angeliki Lazaridou, Quan Ngoc Pham, Raffaella Bernardi, Sandro Pezzelle, Marco Baroni, Gemma Boleda, and Raquel Fernández. The lambada dataset: Word prediction requiring a broad discourse context. \textit{arXiv preprint arXiv:1606.06031}, 2016.}
\refentry{ch08-ref-067}{Kishore Papineni, Salim Roukos, Todd Ward, and Wei-Jing Zhu. Bleu: a method for automatic evaluation of machine translation. In \textit{Proceedings of the 40th Annual} \textit{Meeting on Association for Computational Linguistics}, ACL ’02, page 311–318, USA, 2002. Association for Computational Linguistics. doi: 10.3115/1073083. 1073135. URL \url{https://doi.org/10.3115/1073083.1073135}.}
\refentry{ch08-ref-068}{Arkil Patel, Satwik Bhattamishra, and Navin Goyal. Are nlp models really able to solve simple math word problems? \textit{arXiv preprint arXiv:2103.07191}, 2021.}
\refentry{ch08-ref-069}{Saurav Pawar, S. M Towhidul Islam Tonmoy, S M Mehedi Zaman, Vinija Jain, Aman Chadha, and Amitava Das. The what, why, and how of context length extension techniques in large language models – a detailed survey, 2024.}
\refentry{ch08-ref-070}{Anselmo Peñas, Eduard Hovy, Pamela Forner, Álvaro Rodrigo, Richard Sutcliffe, and Roser Morante. Qa4mre 2011-2013: Overview of question answering for machine reading evaluation. In \textit{Information Access Evaluation. Multilinguality, Multimodality, and Visualization: 4th International Conference of the CLEF Initiative,} \textit{CLEF 2013, Valencia, Spain, September 23-26, 2013. Proceedings 4}, pages 303– 320. Springer, 2013.}
\refentry{ch08-ref-071}{Mohammad Taher Pilehvar and José Camacho-Collados. Wic: 10,000 example pairs for evaluating context-sensitive representations. \textit{arXiv preprint arXiv:1808.09121}, 6:17, 2018.}
\refentry{ch08-ref-072}{Edoardo Maria Ponti, Goran Glavaš, Olga Majewska, Qianchu Liu, Ivan Vulić, and Anna Korhonen. Xcopa: A multilingual dataset for causal commonsense reasoning. \textit{arXiv preprint arXiv:2005.00333}, 2020.}
\refentry{ch08-ref-073}{Jack W Rae, Anna Potapenko, Siddhant M Jayakumar, and Timothy P Lillicrap. Compressive transformers for long-range sequence modelling. \textit{arXiv preprint} \textit{arXiv:1911.05507}, 2019.}
\refentry{ch08-ref-074}{Pranav Rajpurkar, Jian Zhang, Konstantin Lopyrev, and Percy Liang. Squad: 100,000+ questions for machine comprehension of text, 2016.}
\refentry{ch08-ref-075}{Pranav Rajpurkar, Robin Jia, and Percy Liang. Know what you don’t know: Unanswerable questions for SQuAD. In Iryna Gurevych and Yusuke Miyao, editors, \textit{Proceedings of the 56th Annual Meeting of the Association for Computational} \textit{Linguistics (Volume 2: Short Papers)}, pages 784–789, Melbourne, Australia, July 2018. Association for Computational Linguistics. doi: 10.18653/v1/P18-2124. URL \url{https://aclanthology.org/P18-2124}.}
\refentry{ch08-ref-076}{Leonardo Ranaldi, Giulia Pucci, Federico Ranaldi, Elena Sofia Ruzzetti, and Fabio Massimo Zanzotto. Empowering multi-step reasoning across languages via tree-of-thoughts, 2024.}
\refentry{ch08-ref-077}{Hannah Rashkin, Eric Michael Smith, Margaret Li, and Y-Lan Boureau. Towards empathetic open-domain conversation models: A new benchmark and dataset. \textit{arXiv preprint arXiv:1811.00207}, 2018.}
\refentry{ch08-ref-078}{Siva Reddy, Danqi Chen, and Christopher D Manning. Coqa: A conversational question answering challenge. \textit{Transactions of the Association for Computational Linguistics}, 7:249–266, 2019.}
\refentry{ch08-ref-079}{Josh Reini et al. truera/trulens: Trulens eval v0.25.1, 2024. URL \url{https://zenodo.org/doi/10.5281/zenodo.4495856}.}
\refentry{ch08-ref-080}{Melissa Roemmele, Cosmin Adrian Bejan, and Andrew S Gordon. Choice of plausible alternatives: An evaluation of commonsense causal reasoning. In \textit{2011 AAAI} \textit{Spring Symposium Series}, 2011.}
\refentry{ch08-ref-081}{Subhro Roy and Dan Roth. Solving general arithmetic word problems. \textit{arXiv preprint} \textit{arXiv:1608.01413}, 2016.}
\refentry{ch08-ref-082}{Keisuke Sakaguchi, Ronan Le Bras, Chandra Bhagavatula, and Yejin Choi. Winogrande: An adversarial winograd schema challenge at scale. \textit{Communications of} \textit{the ACM}, 64(9):99–106, 2021.}
\refentry{ch08-ref-083}{Maarten Sap, Hannah Rashkin, Derek Chen, Ronan LeBras, and Yejin Choi. Socialiqa: Commonsense reasoning about social interactions. \textit{arXiv preprint} \textit{arXiv:1904.09728}, 2019.}
\refentry{ch08-ref-084}{Elvis Saravia. Prompt Engineering Guide. \textit{https://github.com/dair-ai/Prompt-Engineering-Guide}, 12 2022.}
\refentry{ch08-ref-085}{Roie Schwaber-Cohen. What is a vector database how does it work? use cases + examples, 5 2023. URL \url{https://www.pinecone.io/learn/vector-database/}.}
\refentry{ch08-ref-086}{Thomas Scialom, Paul-Alexis Dray, Sylvain Lamprier, Benjamin Piwowarski, and Jacopo Staiano. Mlsum: The multilingual summarization corpus. \textit{arXiv preprint} \textit{arXiv:2004.14900}, 2020.}
\refentry{ch08-ref-087}{S.D.N.Y. The new york times company v microsoft corporation, openai, inc., openai lp, openai gp, llc, openai llc, openai opco llc, openai global llc, oai corporation, llc and openai holdings llc., 2023. URL \url{https://nytco-assets.nytimes.com/2023/12/NYT_Complaint_Dec2023.pdf}.}
\refentry{ch08-ref-088}{Zhihong Shao, Minlie Huang, Jiangtao Wen, Wenfei Xu, and Xiaoyan Zhu. Long and diverse text generation with planning-based hierarchical variational model. \textit{arXiv preprint arXiv:1908.06605}, 2019.}
\refentry{ch08-ref-089}{Freda Shi, Mirac Suzgun, Markus Freitag, Xuezhi Wang, Suraj Srivats, Soroush Vosoughi, Hyung Won Chung, Yi Tay, Sebastian Ruder, Denny Zhou, et al. Language models are multilingual chain-of-thought reasoners. \textit{arXiv preprint} \textit{arXiv:2210.03057}, 2022.}
\refentry{ch08-ref-090}{Eric Michael Smith, Mary Williamson, Kurt Shuster, Jason Weston, and Y-Lan Boureau. Can you put it all together: Evaluating conversational agents’ ability to blend skills. \textit{arXiv preprint arXiv:2004.08449}, 2020.}
\refentry{ch08-ref-091}{Aarohi Srivastava, Abhinav Rastogi, Abhishek Rao, Abu Awal Md Shoeb, Abubakar Abid, Adam Fisch, Adam R Brown, Adam Santoro, Aditya Gupta, Adrià Garriga-Alonso, et al. Beyond the imitation game: Quantifying and extrapolating the capabilities of language models. \textit{arXiv preprint arXiv:2206.04615}, 2022.}
\refentry{ch08-ref-092}{Alon Talmor, Jonathan Herzig, Nicholas Lourie, and Jonathan Berant. Commonsenseqa: A question answering challenge targeting commonsense knowledge. \textit{arXiv preprint arXiv:1811.00937}, 2018.}
\refentry{ch08-ref-093}{Alexey Tikhonov and Max Ryabinin. It’s all in the heads: Using attention heads as a baseline for cross-lingual transfer in commonsense reasoning. \textit{arXiv preprint} \textit{arXiv:2106.12066}, 2021.}
\refentry{ch08-ref-094}{David Vilares and Carlos Gómez-Rodríguez. Head-qa: A healthcare dataset for complex reasoning. \textit{arXiv preprint arXiv:1906.04701}, 2019.}
\refentry{ch08-ref-095}{Alex Wang, Amanpreet Singh, Julian Michael, Felix Hill, Omer Levy, and Samuel R Bowman. Glue: A multi-task benchmark and analysis platform for natural language understanding. \textit{arXiv preprint arXiv:1804.07461}, 2018.}
\refentry{ch08-ref-096}{Alex Wang, Yada Pruksachatkun, Nikita Nangia, Amanpreet Singh, Julian Michael, Felix Hill, Omer Levy, and Samuel Bowman. Superglue: A stickier benchmark for general-purpose language understanding systems. \textit{Advances in neural information} \textit{processing systems}, 32, 2019.}
\refentry{ch08-ref-097}{Yan Wang, Xiaojiang Liu, and Shuming Shi. Deep neural solver for math word problems. In \textit{Proceedings of the 2017 conference on empirical methods in natural} \textit{language processing}, pages 845–854, 2017.}
\refentry{ch08-ref-098}{Jason Wei et al. Emergent abilities of large language models, 2022.}
\refentry{ch08-ref-099}{Ralph Weischedel, Sameer Pradhan, Lance Ramshaw, Martha Palmer, Nianwen Xue, Mitchell Marcus, Ann Taylor, Craig Greenberg, Eduard Hovy, Robert Belvin, et al. Ontonotes release 4.0. \textit{LDC2011T03, Philadelphia, Penn.: Linguistic Data Consortium}, 17, 2011.}
\refentry{ch08-ref-100}{Johannes Welbl, Nelson F Liu, and Matt Gardner. Crowdsourcing multiple choice science questions. \textit{arXiv preprint arXiv:1707.06209}, 2017.}
\refentry{ch08-ref-101}{WhyLabs. URL \url{https://github.com/whylabs/langkit}.}
\refentry{ch08-ref-102}{Adina Williams, Nikita Nangia, and Samuel R Bowman. A broad-coverage challenge corpus for sentence understanding through inference. \textit{arXiv preprint} \textit{arXiv:1704.05426}, 2017.}
\refentry{ch08-ref-103}{BigScience Workshop et al. Bloom: A 176b-parameter open-access multilingual language model, 2023.}
\refentry{ch08-ref-104}{Yinfei Yang, Yuan Zhang, Chris Tar, and Jason Baldridge. Paws-x: A cross-lingual adversarial dataset for paraphrase identification. \textit{arXiv preprint} \textit{arXiv:1908.11828}, 2019.}
\refentry{ch08-ref-105}{Shukang Yin, Chaoyou Fu, Sirui Zhao, Ke Li, Xing Sun, Tong Xu, and Enhong Chen. A survey on multimodal large language models, 2024.}
\refentry{ch08-ref-106}{Zheng-Xin Yong et al. Prompting multilingual large language models to generate code-mixed texts: The case of south east asian languages, 2023.}
\refentry{ch08-ref-107}{Zhiyuan Yu, Xiaogeng Liu, Shunning Liang, Zach Cameron, Chaowei Xiao, and Ning Zhang. Don’t listen to me: Understanding and exploring jailbreak prompts of large language models, 2024.}
\refentry{ch08-ref-108}{Lifan Yuan, Yangyi Chen, Ganqu Cui, Hongcheng Gao, Fangyuan Zou, Xingyi Cheng, Heng Ji, Zhiyuan Liu, and Maosong Sun. Revisiting out-of-distribution robustness in nlp: Benchmark, analysis, and llms evaluations, 2023.}
\refentry{ch08-ref-109}{Rowan Zellers, Ari Holtzman, Yonatan Bisk, Ali Farhadi, and Yejin Choi. Hellaswag: Can a machine really finish your sentence? \textit{arXiv preprint} \textit{arXiv:1905.07830}, 2019.}
\refentry{ch08-ref-110}{Sheng Zhang, Xiaodong Liu, Jingjing Liu, Jianfeng Gao, Kevin Duh, and Benjamin Van Durme. Record: Bridging the gap between human and machine commonsense reading comprehension. \textit{arXiv preprint arXiv:1810.12885}, 2018.}
\refentry{ch08-ref-111}{Tianyi Zhang, Varsha Kishore, Felix Wu, Kilian Q. Weinberger, and Yoav Artzi. Bertscore: Evaluating text generation with bert, 2020.}
\end{chapterreferences}

